% arara: lualatex
% arara: bib2gls: { group: on }
% arara: lualatex
% arara: bib2gls: { group: on } if found ("log", "Glossary entry `sym.")
% arara: lualatex
% arara: lualatex if found ("log", "Rerun to")
\documentclass[titlepage=false,oneside,
 fontsize=12pt,captions=tableheading]{scrarticle}

\usepackage{pifont}
\usepackage[noatsymgroup]{nlctuserguide}

\renewcommand*{\thispackagename}{glossaries}

\glsxtrnewgls{opt.glostyle.}{\glostyle}
\glsxtrnewgls{opt.acrstyle.}{\acrstyle}
\glsxtrnewgls{opt.abbrstyle.}{\abbrstyle}

\defsemanticcmd{\glostylefmt}{\textsf}{}
\defsemanticcmd{\acrstylefmt}{\textsf}{}
\defsemanticcmd{\abbrstylefmt}{\textsf}{}

\newcommand{\atentry}[1]{\texorpdfstring{\code{@#1}}{@#1}}

\glsxtrnewgls{opt.printgloss.}{\printglossopt}
\newcommand{\printglossoptval}[2]{\optval{printgloss.#1}{#2}}
\newcommand{\printglossoptvalm}[2]{\optval{printgloss.#1}{\marg{#2}}}

\glsxtrnewgls{opt.gloskey.}{\gloskey}
\newcommand{\gloskeyval}[2]{\optval{gloskey.#1}{\marg{#2}}}

\newcommand*{\catfmt}{\csoptfmt}

\glsxtrnewgls{opt.cat.}{\cat}
\glsxtrnewgls{opt.catattr.}{\catattr}
\glsxtrnewgls{opt.resource.}{\resourceopt}
\newcommand{\resourceoptval}[2]{\optval{resource.#1}{#2}}
\newcommand{\resourceoptvalm}[2]{\optval{resource.#1}{\marg{#2}}}

\renewcommand{\optionlistprefix}{idx.opt.}
\renewcommand{\glsxtrtaggedlistsep}{~}
\newcommand{\idxoptiondef}[1]{\hypertarget{option\glsentrytext{idx.opt.#1}}{}\inlineidxdef{opt.#1}}

\renewcommand{\nlctuserguidecustomentryaliases}{%
   glossarystyle=index,
   acronymstyle=index,
   abbreviationstyle=index,
}

\newcommand{\comxr}{\comment{\glslink{pkg.glossaries-extra}{glossaries-extra.sty}}}
\newcommand{\comxronly}{\comment{\glslink{pkg.glossaries-extra}{glossaries-extra.sty} only}}
\newcommand{\comxrkey}{\comment{\glslink{pkg.glossaries-extra}{glossaries-extra.sty} key}}
\newcommand{\comxropt}[1]{\comment{requires \glslink{pkg.glossaries-extra}{glossaries-extra.sty} '\opt{#1}' option}}

\nlctuserguidegls
{
  \def\gprintglossopt#1#2{%
   \glsbibwriteentry{commandoption}{opt.printgloss.#1}{%
    \field{name}{\csoptfmt{#1}}\parent{idx.printglossopt}#2}}%
  \def\gglosty#1#2{\glsbibwriteentry{glossarystyle}{opt.glostyle.#1}%
    {\field{name}{\glostylefmt{#1}}\parent{idx.glossarystyle}#2}}%
  \def\gacrsty#1#2{\glsbibwriteentry{acronymstyle}{opt.acrstyle.#1}%
    {\field{name}{\acrstylefmt{#1}}\parent{idx.acrstyle}#2}}%
  \def\gabbrsty#1#2{\glsbibwriteentry{abbreviationstyle}{opt.abbrstyle.#1}%
    {\field{name}{\abbrstylefmt{#1}}\parent{idx.abbrstyle}#2}}%
  \def\ggloskey#1#2{%
   \glsbibwriteentry{commandoption}{opt.gloskey.#1}{%
    \field{name}{\csoptfmt{#1}}\parent{idx.gloskey}#2}}%
  \def\gglsopt#1#2{%
   \glsbibwriteentry{commandoption}{opt.glsopt.#1}{%
    \field{name}{\csoptfmt{#1}}\parent{gls}#2}}%
  \def\gcat#1#2{%
   \glsbibwriteentry{optionvalue}{opt.cat.#1}{%
    \field{name}{\catfmt{#1}}\parent{idx.category}#2}}%
  \def\gcatattr#1#2{%
   \glsbibwriteentry{optionvalue}{opt.catattr.#1}{%
    \field{name}{\csoptfmt{#1}}\parent{idx.categoryattribute}#2}}%
  \def\gresourceopt#1#2{%
   \glsbibwriteentry{commandoption}{opt.resource.#1}{%
    \field{name}{\csoptfmt{#1}}\parent{GlsXtrLoadResources}#2}}%
  \gidx{category}{\field{plural}{categories}}
  \gidx{glossarystyle}{\name{glossary styles}\field{text}{glossary style}}
  \gidx{acrstyle}{\name{acronym styles}\field{text}{acronym style}
   \field{see}{setacronymstyle}}
  \gidx{abbrstyle}{\name{abbreviation styles}\field{text}{abbreviation style}
   \field{see}{setabbreviationstyle}}
  \gidx{printglossopt}{\name{print [unsrt|noidx] glossary options}%
    \field{text}{print glossary option}%
  }
  % \newglossary
  \gcmd{new\-glossary}
  {
    \providedby{\sty{glossaries}}
    \syntax{\oargm{log-ext}\margm{glossary-label}\margm{out-ext}\margm{in-ext}\margm{title}}
    \desc{defines a \idx{glossary} identified by \meta{glossary-label} (which can
    be referenced by the \gloskey{type} key when defining an entry)}
  }
  % \newglossary*
  \gcmd{new\-glos\-sary*}%
  {% 
    \providedby{\sty{glossaries} v4.08+}
    \syntax{\margm{glossary-label}\margm{title}}
    \desc{a shortcut that supplies file extensions based on the
    \idx{glossary} label:\begin{compactcodebox}%
\gls{newglossary}\oarg{\meta{glossary-label}-glg}\margm{glossary-label}\marg{\meta{glossary-label}-gls}\marg{\meta{glossary-label}-glo}\margm{title}%
\end{compactcodebox}\glsxtrnopostpunc
}
  }
  % \GlsSetQuote
  \gcmd{Gls\-Set\-Quote}
  {
    \providedby{\sty{glossaries} v4.24+}
    \syntax{\margm{character}}
    \desc{set \app{makeindex}['s] quote character to \meta{character}}
  }
  % \glsdisablehyper
  \gcmd{gls\-disable\-hyper}
  {
    \providedby{\sty{glossaries}}%
    \desc{suppresses all \idx{glossary} related hyperlinks}
  }
  % \glsenablehyper
  \gcmd{gls\-enable\-hyper}
  {
    \providedby{\sty{glossaries}}%
    \desc{enables \idx{glossary} related hyperlinks. This is the
    default if \sty{hyperref} is loaded. Has no effect if
    \sty{hyperref} wasn't loaded}
  }
  % \setglossarystyle
  \gcmd{set\-glossary\-style}
  {
    \providedby{\sty{glossaries}}%
    \syntax{\margm{style name}}
    \desc{set the current \idx{glossarystyle} to \meta{style name}}
  }
  % \printglossary
  \gcmd{print\-glos\-sary}
  {
    \providedby{\sty{glossaries}}%
    \syntax{\oargm{options}}
    \desc{displays the \idx{glossary} by inputting a file created by
     \app+{makeindex} or \app+{xindy}. Must be used with
     \gls{makeglossaries} and either \app{makeindex} or \app{xindy}}
  }
  % \printglossaries
  \gcmd{print\-glos\-saries}
  {
    \providedby{\sty{glossaries}}%
    \desc{iterates over all glossaries and does
     \code{\gls{printglossary}\oarg{\printglossoptval{type}{\meta{type}}}}
     for each \idx{glossary}}
  }
  % \printnoidxglossary
  \gcmd{print\-no\-idx\-glos\-sary}%
  {%
    \providedby{\sty{glossaries} v4.04+}
    \syntax{\oargm{options}}
    \desc{displays the \idx{glossary} by obtaining the \gls{indexing} information from
     the \ext+{aux} file and using \TeX\ to sort and collate. Must be used with
     \gls{makenoidxglossaries}. This method can be very slow and has limitations}
  }%
  % \printnoidxglossaries
  \gcmd{print\-no\-idx\-glos\-saries}%
  {%
    \providedby{\sty{glossaries} v4.04+}
    \desc{iterates over all glossaries and does
     \code{\gls{printnoidxglossary}\oarg{\printglossoptval{type}{\meta{type}}}}
     for each \idx{glossary}}
  }
  % \printunsrtglossary
  \gcmd{print\-un\-srt\-glos\-sary}%
  {%
    \providedby{\sty{glossaries-extra} v1.08+}
    \syntax{\oargm{options}}
    \desc{displays the \idx{glossary} by iterating over all entries
     associated with the given \idx{glossary} (in the order in which they
     were added to the \idx{glossary}). The location lists and group headers
     will only be present if the associated fields have been set
     (typically by \app{bib2gls})}
  }%
  % \printunsrtglossaries
  \gcmd{print\-un\-srt\-glos\-saries}%
  {%
    \providedby{\sty{glossaries-extra} v1.08+}
    \desc{iterates over all glossaries and does
     \code{\gls{printunsrtglossary}\oarg{\printglossoptval{type}{\meta{type}}}}
     for each \idx{glossary}}
  }
  % \glsdefaulttype
  \gcmd{gls\-default\-type}%
  {%
    \providedby{\sty{glossaries}}
    \initval{main}%
    \desc{expands to the label of the default \idx{glossary}}
  }%
  % \acronymname
  \gcmd{acronym\-name}%
  {% 
    \providedby{\sty{glossaries}}
    \initval{Acronyms}%
    \note{language-sensitive}
    \desc{expands to the title of the \code{acronym} \idx{glossary}}
  }%
  % \loadglsentries
  \gcmd{load\-gls\-entries}
  {
    \providedby{\sty{glossaries}}%
    \syntax{\oargm{type}\margm{filename}}
    \desc{locally assigns \gls{glsdefaulttype} to \meta{type} and
     inputs \meta{filename}. If the optional argument is omitted,
     the default glossary is assumed. Note that if any entries within
     \meta{filename} have the \gloskey{type} key set (including
     implicitly in commands like \gls{newacronym}), then this will
     override the type given in the optional argument}
  }
  % \makeglossaries
  \gcmd{make\-glos\-saries}
  {
    \providedby{\sty{glossaries}}
    \desc{opens the associated \idx{glossary} files that need to be
     processed by \app+{makeindex} or \app+{xindy}}
  }
  % \makenoidxglossaries
  \gcmd{make\-no\-idx\-glos\-saries}%
  {%
    \providedby{\sty{glossaries} v4.04+}
    \desc{sets up all glossaries so that they can be displayed
    with \gls{printnoidxglossary}}
  }
  % \GlsXtrLoadResources
  \gcmd{Gls\-Xtr\-Load\-Resources}%
  {%
    \providedby{\sty{glossaries-extra} v1.11+}
    \syntax{\oargm{options}}
    \desc{for use with \app{bib2gls}, this both sets up the resource options
     (which \app{bib2gls} can detect from the \ext{aux} file) and
     inputs the \ext{glstex} file created by \app{bib2gls}}
  }
  % \newglossaryentry
  \gcmd{new\-glos\-sary\-entry}
  {
    \providedby{\sty{glossaries}}
    \syntax{\margm{label}\marg{\keyvallist}}
    \desc{defines a new \idx{glossary} entry with the given label. The
    second argument is a comma-separated list of \idxpl{gloskey}}
  }
  % \longnewglossaryentry
  \gcmd{long\-new\-glos\-sary\-entry}%
  {%
    \providedby{\sty{glossaries}}
    \syntax{\margm{label}\margm{key=value list}\margm{description}}
    \desc{defines a new \idx{glossary} entry with the given label. The
    second argument is a comma-separated list of \idxpl{gloskey}.
    The third argument is the description, which may include
    paragraph breaks}
  }
  % \newterm
  \gcmd{new\-term}%
  {%
    \syntax{\oargm{key=value list}\margm{entry-label}}
    \note{requires \opt{index} package option}
    \desc{defines a new \idx{glossary} entry with the given label,
    \gloskey{type} set to \code{index}, the \gloskey{name}
    set to \meta{entry-label} and the \gloskey{description}
    set to \gls{nopostdesc}. The
    optional argument is a comma-separated list of \idxpl{gloskey},
    which can be used to override the defaults}
  }
  % \glsnamefont
  \gcmd{gls\-name\-font}
  {  
    \providedby{\sty{glossaries}}
    \syntax{\margm{text}}
    \desc{used within the predefined glossary styles to apply a font change to the
    \gloskey{name}}
  } 
  % \gls
  \gcmdsp{gls}
  {
    \providedby{\sty{glossaries}}
    \syntax{\oargm{options}\margm{label}\oargm{insert}}
    \desc{references the entry identified by \meta{label}. The text
    produced may vary depending on whether or not this is the \idx{firstuse}}
  }
  % \glspl
  \gcmdsp{glspl}
  {
    \providedby{\sty{glossaries}}
    \syntax{\oargm{options}\margm{label}\oargm{insert}}
    \desc{as \gls{gls} but uses the plural form}
  }
  % \Gls
  \gcmdsp{Gls}
  {
    \providedby{\sty{glossaries}}
    \syntax{\oargm{options}\margm{label}\oargm{insert}}
    \desc{as \gls{gls} but converts the first character of the
    \idx{linktext} to \idx{uppercase}\glsadd{idx.sentencecase} (for
    the start of a sentence) using \gls{makefirstuc}}
  }
  % \Glspl
  \gcmdsp{Glspl}
  {
    \providedby{\sty{glossaries}}
    \syntax{\oargm{options}\margm{label}\oargm{insert}}
    \desc{as \gls{Gls} but uses the plural form}
  }
  % \glssymbol
  \gcmdsp{gls\-symbol}{%
    \syntax{\oargm{options}\margm{label}\oargm{insert}}
    \desc{references the entry identified by \meta{label}. The text
    produced is obtained from the \gloskey{symbol} value.
    The \meta{insert} argument will be inserted at the end of the
    \idx{linktext}. This command does not alter or depend on the
    \idx{firstuseflag}}
  }
  % \glslink
  \gcmdsp{gls\-link}
  {
    \providedby{\sty{glossaries}}
    \syntax{\oargm{options}\margm{label}\margm{text}}
    \desc{references the entry identified by \meta{label} with the
     given \meta{text} as the \idx{linktext}. This command doesn't
     change the \idx{firstuseflag}}
  }
  % \glsadd
  \gcmd{gls\-add}
  {
    \providedby{\sty{glossaries}}
    \syntax{\oargm{options}\margm{label}}
    \desc{indexes the entry identified by \meta{label} without
    producing any text. This ensures the entry appears in the
    \idx{glossary}}
  }
  % \glsfmttext
  \gcmd{gls\-fmt\-text}
  {
     \syntax{\margm{entry-label}}
     \desc{for use within captions or section titles to display the
      formatted \gloskey{text}}
  }
  % \glsfmtshort
  \gcmd{gls\-fmt\-short}
  {
     \syntax{\margm{entry-label}}
     \desc{for use within captions or section titles to display the
      formatted short form}
  }  
  % \glsentrytext
  \gcmd{gls\-entry\-text}
  {%
    \providedby{\sty{glossaries}}
    \syntax{\margm{entry-label}}
    \desc{simply expands to the value of the \gloskey{text} field.
     Does nothing if the entry hasn't been defined. May be used in
     expandable contexts provided that the \gloskey{text} field doesn't
     contain any fragile commands}
  }
  % \glsentryshort
  \gcmd{gls\-entry\-short}
  {%
    \providedby{\sty{glossaries}}
    \syntax{\margm{entry-label}}
    \desc{simply expands to the value of the \gloskey{short} field.
     Does nothing if the entry hasn't been defined. May be used in
     expandable contexts provided that the \gloskey{short} field doesn't
     contain any fragile commands}
  }
  % \glsentrylong
  \gcmd{gls\-entry\-long}
  {%  
    \providedby{\sty{glossaries}}
    \syntax{\margm{entry-label}}
    \desc{simply expands to the value of the \gloskey{long} field.
     Does nothing if the entry hasn't been defined. May be used in
     expandable contexts provided that the \gloskey{long} field doesn't
     contain any fragile commands}
  } 
  % \glsabbrvscfont
  \gcmd{gls\-abbrv\-sc\-font}
  {
    \providedby{\sty{glossaries-extra} v1.17+}
    \syntax{\margm{text}} 
    \desc{short form font used by the small-caps \qt{sc} abbreviation styles}   
  }  
  % \glslongonlyfont
  \gcmd{gls\-long\-only\-font}
  {
    \providedby{\sty{glossaries-extra} v1.17+}
     \syntax{\margm{text}} 
     \desc{long form font used by the \qt{only} abbreviation styles}
  }   
  % \glsabbrvonlyfont
  \gcmd{gls\-abbrv\-only\-font} 
  {
    \providedby{\sty{glossaries-extra} v1.17+}
     \syntax{\margm{text}}
     \desc{short form font used by the \qt{only} abbreviation styles}
  }
  % \glsxtrnewsymbol
  \gcmd{gls\-xtr\-new\-symbol}%
  {%
    \syntax{\oargm{key=value list}\margm{label}\margm{sym}}
    \providedby{\sty{glossaries-extra}}
    \note{requires \code{\csfmt{usepackage}\oarg{\opt{symbols}}\marg{glossaries-extra}}}
    \desc{defines a new \idx{glossary} entry with the given label,
    \gloskey{type} set to \code{symbols}, the \gloskey{category} set
    to \code{symbol}, the \gloskey{name} set to \meta{sym} and the \gloskey{sort}
    set to \meta{label}. The optional argument is a comma-separated list of \idxpl{gloskey},
    which can be used to override the defaults}
  }
  % \newacronym
  \gcmd{new\-acronym}
  {
    \providedby{\sty{glossaries}}
    \syntax{\oargm{options}\margm{label}\margm{short}\margm{long}}
    \desc{this command is provided by the base \sty{glossaries}
     package but is redefined by \sty{glossaries-extra} to use
     \gls{newabbreviation} with the \gloskey{category} key set to
     \cat{acronym}. With just the base \sty{glossaries} package, use
     \gls{setacronymstyle} to set the style. With
     \sty{glossaries-extra}, use
     \code{\gls{setabbreviationstyle}\oarg{\cat{acronym}}\margm{style}} to
     set the style that governs \gls{newacronym}}
  }
  % \setacronymstyle
  \gcmd{set\-acronym\-style}
  {
    \providedby{\sty{glossaries}}
    \syntax{\oargm{glossary-type}\margm{style-name}}
    \desc{sets the acronym style. Don't use with \sty{glossaries-extra}}
  }
  % \newabbreviation
  \gcmd{new\-abbreviation}
  {
    \providedby{\sty{glossaries-extra}}
    \syntax{\oargm{options}\margm{label}\margm{short}\margm{long}}
    \desc{defines a new entry that represents an abbreviation.
     This internally uses \gls{newglossaryentry} and any provided
     \meta{options} (\idxpl{gloskey}) will be appended. The
     \gloskey{category} is set to \cat{abbreviation} by default, but
     may be overridden in \meta{options}. The
     appropriate style should be set before the abbreviation is
     defined with \gls{setabbreviationstyle}}
  }
  % \setabbreviationstyle
  \gcmd{set\-abbreviation\-style}
  {
    \providedby{\sty{glossaries-extra}}
    \syntax{\oargm{category}\margm{style-name}}
     \desc{sets the current \idx{abbrstyle} to \meta{style-name}
      for the category identified by \meta{category}. If the
      optional argument is omitted, \cat{abbreviation} is assumed}
  }
  % \acrlong
  \gcmdsp{acr\-long}
  {
    \providedby{\sty{glossaries}}
    \syntax{\oargm{options}\margm{label}\oargm{insert}}
    \desc{displays the long form of an acronym. Only for use with
     the base \sty{glossaries} package's acronym mechanism. With
    \sty{glossaries-extra}, use \gls{glsxtrlong} instead}
  }
  % \acrshort
  \gcmdsp{acr\-short}
  {
    \providedby{\sty{glossaries}}
    \syntax{\oargm{options}\margm{label}\oargm{insert}}
    \desc{displays the short form of an acronym. Only for use with
     the base \sty{glossaries} package's acronym mechanism. With
    \sty{glossaries-extra}, use \gls{glsxtrshort} instead}
  }
  % \glsxtrlong
  \gcmdsp{gls\-xtr\-long}
  {
    \providedby{\sty{glossaries-extra}}
    \syntax{\oargm{options}\margm{label}\oargm{insert}}
    \desc{displays the long form of an abbreviation}
  }
  % \glsxtrshort
  \gcmdsp{gls\-xtr\-short}
  {
    \providedby{\sty{glossaries-extra}}
    \syntax{\oargm{options}\margm{label}\oargm{insert}}
    \desc{displays the short form of an abbreviation}
  }
  % \ifglsused
  \gcmd{if\-gls\-used}
  {
    \providedby{\sty{glossaries}}
    \syntax{\margm{label}\margm{true text}\margm{false text}}
    \desc{does \meta{true} if the entry has been marked as
    \idxc{firstuseflag}{used} and does \meta{false} if the entry is
    marked as \idxc{firstuseflag}{unused}}
  }
  % \glsnoexpandfields
  \gcmd{gls\-no\-expand\-fields}
  {
    \providedby{\sty{glossaries}}
    \desc{don't expand field values when defining entries, except
     for those that explicitly have expansion enabled}
  }
  % \glsreset
  \gcmd{gls\-reset}
  {
    \providedby{\sty{glossaries}}
    \syntax{\margm{label}}
    \desc{globally resets the entry's \idx{firstuseflag}. That is, this marks the entry
     as \qt{not used}}
  }
  % \glsunset
  \gcmd{gls\-unset}
  {
    \providedby{\sty{glossaries}}
    \syntax{\margm{label}}
    \desc{globally unsets the entry's \idx{firstuseflag}. That is, this marks the entry
     as \qt{used}}
  }
  % \glsresetall
  \gcmd{gls\-reset\-all}
  {
    \providedby{\sty{glossaries}}
    \syntax{\oargm{types}}
    \desc{globally resets all entries associated with the listed
     glossaries or all glossaries if \meta{types} is omitted}
  }
  % \glsunsetall
  \gcmd{gls\-unset\-all}
  {
    \providedby{\sty{glossaries}}
    \syntax{\oargm{types}}
    \desc{globally unsets all entries associated with the listed
     glossaries or all glossaries if \meta{types} is omitted}
  }
  % \glslowercase
  \gcmd{gls\-lower\-case}
  {
    \providedby{\sty{glossaries} v4.50+}
    \syntax{\margm{text}}
    \desc{converts \meta{text} to \idx{lowercase} using the modern \LaTeX3
    case-changing command, which is expandable}
  }
  % \makefirstuc
  \gcmd{make\-first\-uc}%
  {%
    \syntax{\margm{text}}
    \providedby{\sty{mfirstuc}}
    \desc{robust command that converts the first character of \meta{text} to
    \idx{uppercase} (\idx{sentencecase}). See the \sty{mfirstuc} 
    documentation for further details, either:
    \texdocref{mfirstuc} or visit \ctanpkg{mfirstuc}}
  }
  % \nopostdesc
  \gcmd{no\-post\-desc}
  {%
    \providedby{\sty{glossaries} v1.17+}
    \desc{when placed at the end of the \gloskey{description}, this
     switches off the post-description punctuation (if it has been
     enabled). Does nothing outside of the \idx{glossary}}
  }
  % \seename
  \gcmd{see\-name}
  {
    \providedby{\sty{glossaries}}
    \initval{see}
    \note{language-sensitive}
    \desc{used as a cross-reference tag}
  }
  % \seealsoname
  \gcmd{see\-also\-name}
  {
    \providedby{\sty{glossaries-extra}}
    \initval{see also}
    \note{language-sensitive}
    \desc{used as a cross-reference tag}
  }
  % \newglossaryentry (glossentry) keys
  \gidxpl{gloskey}{%
    \common
    \field{text}{glossary entry key}
    \desc{these are options that can be passed to commands that
     define entries, such as \gls{newglossaryentry} or
     \gls{newabbreviation}}
  }
  % glossentry name
  \ggloskey{name}%
  {%
    \syntax{\margm{text}}
    \desc{the entry's name, as displayed in the \idx{glossary}}
  }
  % glossentry description
  \ggloskey{description}%
  {%
    \syntax{\margm{text}}
    \desc{the entry's description, as displayed in the \idx{glossary}}
  }
  % glossentry type
  \ggloskey{type}%
  {%
    \syntax{\meta{\idx{glossary}-label}}
    \initval{\gls{glsdefaulttype}}%
    \desc{assigns the entry to the \idx{glossary} identified by
      \meta{\idx{glossary}-label}}
  }
  % glossentry parent
  \ggloskey{parent}%
  {%
    \syntax{\meta{parent-label}}
    \desc{the label of the entry's parent (from which the entry's
    hierarchical level is obtained)}
  }
  % glossentry category
  \ggloskey{category}%
  {%
    \syntax{\meta{category-label}}%
    \initval{general}%
    \providedby{\sty{glossaries-extra}}
    \desc{the entry's \idx{category} (must be a simple label)}
  }
  % glossentry sort
  \ggloskey{sort}%
  {%
    \syntax{\meta{value}}
    \initval{\meta{entry name}}
    \desc{specifies the value to use for sorting (overrides the
     default). This key is usually required for \app+{xindy} if the
      \gloskey{name} key only contains commands (for example, the
      entry is a symbol), but explicitly using this key in other
      contexts can break certain sort methods.
      \gallerypage{bib2gls-sorting}{Don't use the \gloskey{sort} field
      with \app{bib2gls}}}
  }
  % glossentry see
  \ggloskey{see}%
  {%
    \syntax{\oargm{tag}\margm{label-list}}
    \desc{indexes each item in \meta{label-list} with the
     cross-reference in place of a normal location}
  }
  % glossentry seealso
  \ggloskey{see\-also}%
  {%
    \providedby{\sty{glossaries-extra}}
    \syntax{\margm{label-list}}
    \desc{like \gloskey{see} but always uses \gls{seealsoname} as the tag}
  }
  % glossentry alias
  \ggloskey{alias}%
  {%
    \providedby{\sty{glossaries-extra}}
    \syntax{\margm{xr-label}}
    \desc{like \gloskey{see} but only allows a single
    cross-reference with no optional tag, and sets up aliasing, which
    changes the hyperlink target}
  }
  % glossentry text
  \ggloskey{text}%
  {%
    \syntax{\margm{text}}
    \desc{the entry's text, as displayed on \idx{subsequentuse} of
     \gls{glslike} commands. If omitted, this value is assumed to be
     the same as the \gloskey{name} key}
  }
  % glossentry plural
  \ggloskey{plural}%
  {%
    \syntax{\margm{text}}
    \desc{the entry's plural form, as displayed on \idx{subsequentuse} of
     plural \gls{glslike} commands, such as \gls{glspl}. This should
     be the appropriate plural form of the value provided by the
    \gloskey{text} key. If omitted, this value is assumed to be the
    value of the \gloskey{text} key with \gls{glspluralsuffix}
    appended}
  }
  % glossentry symbol
  \ggloskey{symbol}%
  {%
    \initval{\gls{relax}}
    \syntax{\margm{symbol}}
    \desc{the entry's associated symbol (optional), which can be
     displayed with \gls{glssymbol} (if \gls{indexing} and hyperlinks are
     required) or with \gls{glsentrysymbol}}
  }
  % glossentry first
  \ggloskey{first}%
  {%
    \syntax{\margm{first}}
    \desc{the entry's text, as displayed on \idx{firstuse} of
     \gls{glslike} commands. If omitted,
     this value is assumed to be the same as the \gloskey{text} key}
  }
  % glossentry firstplural
  \ggloskey{first\-plural}%
  {%
    \syntax{\margm{text}}
    \desc{the entry's plural form, as displayed on \idx{firstuse} of
     plural \gls{glslike} commands, such as \gls{glspl}. If this key
    is omitted, then the value will either be the same as the
    \gloskey{plural} field, if the \gloskey{first} key wasn't used, or
    the value will be taken from the \gloskey{first} key with
    \gls{glspluralsuffix} appended}
  }
  % glossentry short
  \ggloskey{short}%
  {%
    \syntax{\margm{short form}}
    \desc{a key that is set implicitly by \gls{newabbreviation}
     and \gls{newacronym} to the entry's short (abbreviated) form. It
     typically shouldn't be used explicitly with
     \gls{newglossaryentry}. However, if you are using \app{bib2gls}
     then this field should be used in the \ext{bib} file when
     defining abbreviations}
  }
  % glossentry long
  \ggloskey{long}%
  {%
    \syntax{\margm{long form}}
    \desc{a key that is set implicitly by \gls{newabbreviation}
     and \gls{newacronym} to the entry's long (unabbreviated) form. It
     typically shouldn't be used explicitly with
     \gls{newglossaryentry}. However, if you are using \app{bib2gls}
     then this field should be used in the \ext{bib} file when
     defining abbreviations}
  }
  % glossentry shortplural
  \ggloskey{short\-plural}%
  {%
    \syntax{\margm{plural short form}}
    \desc{the plural corresponding to the \gloskey{short} key if the
    default value is inappropriate}
  }
  % glossentry longplural
  \ggloskey{long\-plural}%
  {%
    \syntax{\margm{plural long form}}
    \desc{the plural corresponding to the \gloskey{long} key if the
    default value is inappropriate}
  }
  % glossentry user1
  \ggloskey{user1}%
  {%
    \syntax{\margm{text}}
    \desc{a generic field, which can be
     displayed with \gls{glsuseri} (if \gls{indexing} and hyperlinks are
     required) or with \gls{glsentryuseri}}
  }
  % package option style
  \gstyopt{style}%
  {%
    \syntax{\meta{style-name}}
  }
  % package option nogroupskip
  \gstyopt{no\-group\-skip}
  {
    \initval{false}
    \defval{true}
    \syntax{\meta{boolean}}
  }
  % package option nopostdot
  \gstyopt{no\-post\-dot}
  {
    \initval{false}
    \defval{true}
    \syntax{\meta{boolean}}
  }
  % package option postdot
  \gstyopt{post\-dot}{\providedby{\sty{glossaries-extra}}}
  % package option nonumberlist
  \gstyopt{no\-number\-list}{}
  % package option symbols
  \gstyopt{symbols}{}
  % package option numbers
  \gstyopt{numbers}{}
  % package option index
  \gstyopt{index}{}
  % package option abbreviations
  \gstyopt{abbreviations}{}
  % package option counter
  \gstyopt{counter}%
  {%
    \syntax{\meta{counter-name}}
  }
  % package option toc
  \gstyopt{toc}%
  {%
    \syntax{\meta{boolean}}
  }
  % package option section
  \gstyopt{section}
  {
    \syntax{\meta{value}}
  }
  % package option acronym
  \gstyopt{acronym}%
  {%
    \initval{false}
    \defval{true}
    \syntax{\meta{boolean}}
  }
  % package option acronyms
  \gstyopt{acronyms}{}
  % package option stylemods
  \gstyopt{style\-mods}{}
  % package option sort
  \gstyopt{sort}%
  {%
    \syntax{\meta{setting}}
  }
  % package option sanitizesort
  \gstyopt{sanitize\-sort}{}
  % package option record
  \gstyopt{record}%
  {%
    \initval{off}
    \defval{only}
    \syntax{\meta{setting}}
  }
  % package option automake
  \gstyopt{auto\-make}{}
  % package option xindy
  \gstyopt{xindy}{}
  % package option order
  \gstyopt{order}{}
  % package option nohypertypes
  \gstyopt{no\-hyper\-types}{\syntax{\margm{label list}}}
  % package option hyperfirst
  \gstyopt{hyper\-first}
  {
     \initval{true}
     \defval{true}
     \syntax{\meta{boolean}}
  }
  % category: abbreviation
  \gcat{abbreviation}{}
  % category: acronym
  \gcat{acronym}{}
  % glossary style: index
  \gglosty{index}{}
  % glossary style: indexgroup
  \gglosty{index\-group}{}
  % glossary style: mcolindex
  \gglosty{mcol\-index}{}
  % glossary style: altlist
  \gglosty{alt\-list}{}
  % glossary style: long4col
  \gglosty{long4col}{}
  % acronym style: long-short
  \gacrsty{long\dhyphen short}{}
  % acronym style: long-sc-short
  \gacrsty{long\dhyphen sc\dhyphen short}{}
  % acronym style: long-short-desc
  \gacrsty{long\dhyphen short\dhyphen desc}{}
  % acronym style: short-long
  \gacrsty{short\dhyphen long}{}
  % acronym style: short-long-desc
  \gacrsty{short\dhyphen long\dhyphen desc}{}
  % abbreviation style: long-short
  \gabbrsty{long\dhyphen short}{}
  % abbreviation style: long-short-sc
  \gabbrsty{long\dhyphen short\dhyphen sc}{}
  % abbreviation style: short-nolong
  \gabbrsty{short\dhyphen nolong}{}
  % abbreviation style: long-only-short-only
  \gabbrsty{long\dhyphen only\dhyphen short\dhyphen only}{}
  % resource option src
  \gresourceopt{src}{\syntax{bib list}}
  % resource option type
  \gresourceopt{type}{\syntax{glossary-label}}
  % resource option sort
  \gresourceopt{sort}{\syntax{method}}
  % resource option sort-field
  \gresourceopt{sort\dhyphen field}{\syntax{field name}}
  % resource option save-locations
  \gresourceopt{save\dhyphen locations}{\syntax{\meta{boolean}}}
  % printgloss style
  \gprintglossopt{style}{%
    \syntax{\meta{style-name}}
  }
  % printgloss type
  \gprintglossopt{type}{%
    \syntax{\meta{glossary-label}}
  }
  % printgloss nonumberlist
  \gprintglossopt{no\-number\-list}{}
  % printgloss toctitle
  \gprintglossopt{toc\-title}{}
  % printgloss sort
  \gprintglossopt{sort}{}
  % general commands and environments
  \gcmd{textsc}{}% \textsc
  \gcmd{item}{}% \item
  \gcmd{tex\-or\-pdf\-string}{}% \texorpdfstring
  \genv{description}{}
  % applications:
  \gapp{make\-index}{}
  \gapp{xindy}{}
  \gapp{make\-glos\-saries}{}
  \gapp{make\-glos\-saries\dhyphen lite}{}% makeglossaries-lite
  \gapp{bib2gls}{}
  \gapp{arara}{}
  % --group
  \glongswitch{group}{\inapp{bib2gls}}
  % -g
  \gshortswitch{g}{\inapp{bib2gls}\field{alias}{switch.group}}
  % packages
  \gpkg{glossaries}{\common}% glossaries.sty
  \gpkg{glossaries\dhyphen extra}{\common}% glossaries-extra.sty
  \gpkg{glossary\dhyphen mcols}{}% glossary-mcols.sty
  \gpkg{data\-tool}{}% datatool.sty
  \gpkg{data\-tool\dhyphen base}{}% datatool-base.sty
  \gpkg{mfirst\-uc}{}% mfirstuc.sty
  \gpkg{hyper\-ref}{}% hyperref.sty
  \gpkg{babel}{}% babel.sty
  \gpkg{poly\-glossia}{}% polyglossia.sty
  \gpkg{input\-enc}{}% inputenc.sty
  \gpkg{font\-enc}{}% fontenc.sty
  % files
  \gfilemeta{glossaries\dhyphen}{language}{}{}% glossaries-<language>
  \gfile{glossaries\dhyphen french}{}% glossaries-french
  \gfile{glossaries\dhyphen german}{}% glossaries-german
  \gext{log}{}
  \gext{tex}{}
  \gext{bib}{}
  \gext{glo}{}
  \gext{gls}{}
  \gext{glg}{}
  \gext{ist}{}
  \gext{xdy}{}
  \gext{aux}{}
  \gext{glstex}{}
  \gext{alg}{}
  \gext{acr}{}
  \gext{acn}{}
  % terms
  \gterm{field}%
  {%
    \name{field}
    \desc{entry data is stored in fields. These may have a
     corresponding key used to set the value, such as \gloskey{name}
     or \gloskey{description}}
  }
  \gterm{indexing}{%
    \name{indexing (or recording)}
    \field{text}{indexing}
    \desc{the process of saving the location (for example, the page number) and any
      associated information that is required in the \idx{glossary}.
      The data may be sorted and collated by an indexing application
      as part of the document build}
  }
  \gterm{glossary}%
  {%
    \common
    \name{glossary}
    \field{plural}{glossaries}
    \desc{technically a glossary is an alphabetical list of words
    relating to a particular topic. For the purposes of describing
    the \styfmt{glossaries} and \sty{glossaries-extra} packages, a
    glossary is either the list produced by commands like
    \gls{printglossary} or \gls{printunsrtglossary} (which may or
    may not be ordered alphabetically) or a glossary is a set of
    entry labels where the set is identified by the glossary label
    or type}
  }
  \gterm{firstuseflag}{\name{first use flag}%
    \common
    \desc{a conditional that keeps track of whether or not an entry
     has been referenced by any of the \gls{glslike} commands (which
     can adjust their behaviour according to whether or not this flag is
     true). The conditional is true if the entry hasn't been used by
     one of these commands
     (or if the flag has been reset) and false if it has been used
     (or if the flag has been unset)}
  }
  \gterm{firstuse}{\name{first use}%
    \common
    \desc{the first time an entry is used by a command that unsets
     the \idx{firstuseflag} (or the first time since the flag was
     reset)}
  }
  \gterm{firstusetext}{\name{first use text}%
    \common
    \desc{the \idx{linktext} that is displayed on \idx{firstuse} of
    the \gls{glslike} commands}
  }
  \gterm{subsequentuse}{\name{subsequent use}
    \common
    \desc{using an entry that unsets the \idx{firstuseflag} when it
    has already been unset}
  }
  \gterm{glslike}{\common\name{{}\csfmt{gls}-like}
    \desc{commands like \gls{gls} that change the
    \idx{firstuseflag}. These commands index the entry (if indexing
    is enabled), create a hyperlink to the entry's \idx{glossary} listing (if enabled)
    and unset the \idx{firstuseflag}}
  }
  \gterm{linktext}{\name{link text}%
    \desc{the text produced by commands like \gls{gls} that have the 
    potential to be a hyperlink}
  }
  \gterm{semanticcommand}{\name{semantic command}
    \desc{essentially, a semantic command is one that's associated with a
    particular document element, idea or topic that hides the font
    and other stylistic formatting inside its definition. For
    example, Latin taxonomy is usually displayed in italic.
    Explicitly using font commands, for example
    \code{\cmd{textit}\marg{Clostridium}}, is syntactic markup.
    Whereas defining a command called, say, \csfmt{bacteria} that
    displays its argument in italics is a semantic command. The
    actual styling is hidden and the command relates specifically to
    a particular concept.
\begin{codebox}
Syntactic: \cmd{textit}\marg{Clostridium}\codepar
\cmd{newrobustcmd}*\marg{\cmd{bacteria}}[1]\marg{\cmd{emph}\marg{\#1}}\%\newline
Semantic: \cmd{bacteria}\marg{Clostridium}
\end{codebox}
    The end result is the same:
\begin{resultbox}
Syntactic: \textit{Clostridium}\glspar
Semantic: \emph{Clostridium}
\end{resultbox}
    The advantage with semantic commands is that it's much easier to
    change the style, simply by adjusting the command definition.
    Note that I've used \csfmt{newrobustcmd} to make the semantic
    command robust as the style commands can cause a problem if they
    expand too soon. Alternatively, use \csfmt{NewDocumentCommand}}
  }
  % index only
  \gidx{uppercase}{\field{seealso}{idx.sentencecase}}
  \gidx{lowercase}{}
  \gidx{sentencecase}{\name{sentence case}}
  \gidx{allcaps}{\name{all caps}}
  % punctuation
  \gpunc{sym.nbsp}{\name{\code{\textasciitilde} (non-breaking space)}
    \field{see}{idx.nbsp}
  }
  \gidx{nbsp}{\name{non-breaking space (\code{\textasciitilde})}
    \field{symbol}{\code{\textasciitilde}}
  }
  \gpunc{sym.colon}{\name{\code{:} (colon)}}
  \gpunc{unknowntag}{\name{\code{??} (unknown entry marker)}
    \field{text}{\code{??}}}
  % indexing options
  \gidx{opt.noidx}{\name{Option 1 (\qt{noidx})}\field{text}{1}}
  \gidx{opt.mkidx}{\name{Option 2 (\appfmt{makeindex})}\field{text}{2}}
  \gidx{opt.xdy}{\name{Option 3 (\appfmt{xindy})}\field{text}{3}}
  \gidx{opt.b2g}{\name{Option 4 (\appfmt{bib2gls})}\field{text}{4}}
}

\title{The glossaries package v4.56: 
a guide for beginners}
\author{Nicola L.C. Talbot\\\href{https://www.dickimaw-books.com/}{\nolinkurl{dickimaw-books.com}}}
\date{2025-03-03}

\begin{document}
\maketitle

\begin{abstract}
The \sty{glossaries} package is very flexible, but this means
that it has a~lot of options, and since a user guide is supposed to
provide a complete list of all the high-level user commands, the main
user manual is quite big. This can be rather
daunting for beginners, so this document is a brief introduction
just to help get you started. If you find yourself saying, 
\qt{but how can I do\ldots?} then it's time to move on to the 
\docref{main user manual}{glossaries-user}.

I've made some statements in this document that don't actually tell
you the full truth, but it would clutter the document and cause
confusion if I keep writing \qt{except when \ldots} or \qt{but you can
also do this, that or the other} or \qt{you can do it this way but
you can also do it that way, but that way may cause complications
under certain circumstances}.
\end{abstract}

\tableofcontents

\section{Getting Started}
\label{sec:start}

As with all packages, you need to load \sty+{glossaries} with
\csfmt{usepackage}, but there are certain packages that must be loaded
before \styfmt{glossaries}, \emph{if} they are required: \sty{hyperref},
\sty{babel}, \sty{polyglossia}, \sty{inputenc} and \sty{fontenc}.
(You don't have to load these packages, but if you want them, you
must load them before \sty{glossaries}.)

\begin{important}
If you require multilingual support you must also install 
the relevant language module. Each language module is called
\file{glossaries-language}, where \meta{language} is the
root language name. For example, \file{glossaries-french}
or \file{glossaries-german}. If a language module is required,
the \sty{glossaries} package will automatically try to load it
and will give a warning if the module isn't found.
\end{important}

Once you have loaded the \sty{glossaries} package, you need to define
your terms in the preamble and then you can use them throughout the
document. Here's a simple example:
\begin{codebox*}
\cmd{documentclass}\marg{article}
\cmd{usepackage}\marg{glossaries}
\comment{define a term:}
\gls{newglossaryentry}\marg{\strong{ex}}\marg{\gloskeyval{name}{sample},
 \gloskeyval{description}{an example}}
\cbeg{document}
Here's my \gls{gls}\marg{\strong{ex}} term.
\cend{document}
\end{codebox*}
The above defines a simple term identified by the label \qt{ex}.
When the term is referenced by its label with \gls{gls}, the
provided text (\qt{sample}) is inserted into the document:
\begin{resultbox}
Here's my sample term.
\end{resultbox}
Note that a description has been given in the definition but it doesn't
appear in the document yet.

In general, it's best to choose a label that's easy to remember and
makes the source code (reasonably) easy to read. In the above example, a better
label would be \qt{sample}, the same as the term itself. However,
bear in mind that labels mustn't contain non-expandable content.
The next example defines a term that's actually a symbol. 
In this case, the label must be
different from the provided term, which contains the commands to
typeset the symbol:
\begin{codebox*}
\cmd{documentclass}\marg{article}
\cmd{usepackage}\marg{glossaries}
\gls{newglossaryentry}\marg{\strong{emptyset}}\marg{
 \gloskeyval{name}{\cmd{ensuremath}\marg{\cmd{emptyset}}},
 \gloskeyval{description}{the empty set}}
\cbeg{document}
The empty set is denoted \gls{gls}\marg{\strong{emptyset}}.
\cend{document}
\end{codebox*}
This produces:
\begin{resultbox}
The empty set is denoted $\emptyset$.
\end{resultbox}

The next example defines a term with the label \qt{svm} that has both 
a long form (\qt{support vector machine}) and a short form (\qt{SVM}):
\begin{codebox*}
\cmd{documentclass}\marg{article}
\cmd{usepackage}\marg{glossaries}
\gls{setacronymstyle}\marg{\acrstyle{long-short}}
\gls{newacronym}\marg{\strong{svm}}\marg{SVM}\marg{support vector machine}
\cbeg{document}
First use: \gls{gls}\marg{\strong{svm}}. Second use: \gls{gls}\marg{\strong{svm}}.
\cend{document}
\end{codebox*}
This has used a different command to define the term, but the term is
still referenced with \gls{gls} in the same way as before. However,
the text produced by \gls{gls} now varies:
\begin{resultbox}
First use: support vector machine (SVM). Second use: SVM.
\end{resultbox}
In this case, the text (the \qt{\idx{linktext}}) produced by 
\code{\gls{gls}\marg{svm}}
changed after the \idx{firstuse}. The \idx{firstuse} produced the
long form followed by the short form in parentheses because I set
the acronym style to \acrstyle{long-short}. The \idx{subsequentuse}
just showed the short form.

I suggest you try the above examples to make sure you have the
package correctly installed.  If you get an \qt{undefined control
sequence} error for any of the examples provided in this guide, check that the version
number at the top of this document matches the version you have
installed. (Open the \ext+{log} file and search for the line that
starts with \qt{Package: glossaries} followed by a date and
version.) 

\begin{important}
Be careful of fragile commands. If a command causes a problem when
used in one of the \gls{newglossaryentry} fields, consider adding
\gls{glsnoexpandfields} before you start defining your entries.
Where possible use robust \idxpl{semanticcommand}.
\end{important}

Abbreviations are slightly different if you use the extension package
\sty+{glossaries-extra} (which needs to be installed separately).
This automatically loads \sty{glossaries} so you don't need to
explicitly load it:
\begin{codebox*}
\cmd{documentclass}\marg{article}
\cmd{usepackage}\marg{glossaries-extra}
\comment{commands provided by \sty{glossaries-extra}:}
\gls{setabbreviationstyle}\marg{\abbrstyle{long-short}}
\gls{newabbreviation}\marg{svm}\marg{SVM}\marg{support vector machine}
\cbeg{document}
First use: \gls{gls}\marg{svm}. Second use: \gls{gls}\marg{svm}.
\cend{document}
\end{codebox*}
Since \abbrstyle{long-short} happens to be the default for
\gls{newabbreviation} you may omit the \gls{setabbreviationstyle} 
line in this example.

If you still want to use \gls{newacronym} (rather than
\gls{newabbreviation}) then you need the optional
argument of \gls{setabbreviationstyle}:
\begin{codebox*}
\cmd{documentclass}\marg{article}
\cmd{usepackage}\marg{glossaries-extra}
\gls{setabbreviationstyle}\oarg{\cat{acronym}}\marg{\abbrstyle{long-short}}
\gls{newacronym}\marg{svm}\marg{SVM}\marg{support vector machine}
\cbeg{document}
First use: \gls{gls}\marg{svm}. Second use: \gls{gls}\marg{svm}.
\cend{document}
\end{codebox*}
In this example, if you omit the \gls{setabbreviationstyle} line you
will notice a difference because the \abbrstyle{short-nolong} style (not
the \abbrstyle{long-short} style) is the default with \gls{newacronym}.
With the \abbrstyle{short-nolong} style the first use simply shows just
the short form.

\begin{important}
You can't use \gls{setacronymstyle} with \sty{glossaries-extra}.
\end{important}

If you like, you can put all your definitions in another file (for
example, \filefmt{defns.tex}) and load that file in the preamble using
\gls{loadglsentries} with the filename as the argument. For example:
\begin{codebox}
\gls{loadglsentries}\marg{defns}
\end{codebox}
If you find you have a really large number of definitions that are
hard to manage in a \ext+{tex} file, you might want to have a
look at \app{bib2gls} (installed separately) which requires 
a \ext+{bib} format instead.

Avoid explicitly using formatting commands in abbreviation definitions 
as they can interfere with the underlying mechanism. Instead, the
formatting should be done by the style. For example, suppose I~want
to replace \code{SVM} with \code{\gls{textsc}\marg{svm}}, then
I~need to select a style that uses \gls{textsc}, like this (for the
base \sty{glossaries} style):
\begin{codebox*}
\cmd{documentclass}\marg{article}
\cmd{usepackage}\marg{glossaries}
\gls{setacronymstyle}\marg{\strong{\acrstyle{long-sc-short}}}
\gls{newacronym}\marg{svm}\marg{svm}\marg{support vector machine}
\cbeg{document}
First use: \gls{gls}\marg{svm}. Second use: \gls{gls}\marg{svm}.
\cend{document}
\end{codebox*}
The abbreviation styles have a different naming scheme with
\sty{glossaries-extra}:
\begin{codebox*}
\cmd{documentclass}\marg{article}
\cmd{usepackage}\marg{glossaries-extra}
\comment{commands provided by \sty{glossaries-extra}:}
\gls{setabbreviationstyle}\marg{\strong{\abbrstyle{long-short-sc}}}
\gls{newabbreviation}\marg{svm}\marg{svm}\marg{support vector machine}
\cbeg{document}
First use: \gls{gls}\marg{svm}. Second use: \gls{gls}\marg{svm}.
\cend{document}
\end{codebox*}
With \sty{glossaries-extra} you can have different abbreviation
styles for different categories. Many of these styles have their
own associated formatting commands that can be redefined 
for minor adjustments. See the \sty{glossaries-extra} manual for
further details or \gallerypage{sample-name-font}{Gallery: Mixing Styles}.

As you can hopefully see from the above examples, there are two main ways of
defining a term: as a general entry (\gls{newglossaryentry}) or as an
abbreviation (\gls{newacronym} or, with \sty{glossaries-extra},
\gls{newabbreviation}).

\begin{information}
Regardless of the method of defining a term, the term is always 
given a label which is used to uniquely identify the term (like the standard
\csfmt{label}\slash\csfmt{ref} or \csfmt{cite} mechanism). The label may be
the same as the text produced with \gls{gls} (provided it doesn't
contain any formatting commands) or may be completely different.
\end{information}

The labels are identified in bold in the following:
\begin{codebox}
\gls{newglossaryentry}\marg{\strong{elite}}\marg{\gloskeyval{name}{\'elite},
 \gloskeyval{description}{select group}}
\gls{newglossaryentry}\marg{\strong{set}}\marg{\gloskeyval{name}{set},
 \gloskeyval{description}{collection of distinct elements}}
\gls{newglossaryentry}\marg{\strong{sym.set}}\marg{
 \gloskeyval{name}{\cmd{ensuremath}\marg{\cmd{mathcal}\marg{S}}},
 \gloskeyval{description}{a set}}
\gls{newacronym}\marg{\strong{tool.cnc}}\marg{CNC}\marg{computer numerical control}
\gls{newacronym}\marg{\strong{police.cnc}}\marg{CNC}\marg{civil nuclear constabulary}
\gls{newacronym}\marg{\strong{miltary.cnc}}\marg{CNC}\marg{commander in chief}
\end{codebox}

With modern \TeX\ installations you should now be able to use UTF-8
characters in the label, but beware of active characters. For
example, \sty{babel} makes some punctuation characters, such as
\idx{sym.colon}, active for certain languages. This means that the
character behaves like a command which is unsuitable for a label.

For example, the following works:
\begin{codebox}
\cmd{documentclass}\marg{article}
\cmd{usepackage}\marg{glossaries}
\gls{newglossaryentry}\marg{sym:set}\marg{\gloskeyval{name}{\cmd{ensuremath}\marg{\cmd{mathcal}\marg{S}}},
 \gloskeyval{description}{a set}}
\cbeg{document}
\gls{gls}\marg{sym:set}
\cend{document}
\end{codebox}
However, if \sty{babel} is loaded with \optfmt{french} then the
\idx{sym.colon} character becomes active.
\begin{badcodebox}
\cmd{documentclass}\marg{article}
\cmd{usepackage}[T1]\marg{fontenc}
\cmd{usepackage}[french]\marg{babel}
\cmd{usepackage}\marg{glossaries}
\codepar
\comment{the colon : is a normal character here}
\gls{newglossaryentry}\marg{sym:set}\marg{
 \gloskeyval{name}{\cmd{ensuremath}\marg{\cmd{mathcal}\marg{S}}},
 \gloskeyval{description}{a set}}
\codepar
\cbeg{document}
\comment{the colon : is now an active character}
\gls{gls}\marg{sym:set}
\cend{document}
\end{badcodebox}
You may find that the above example seems to work, but a problem
will occur if \sty{hyperref} and a \idx{glossary} list are added to the
document as the active character will interfere with the hyperlink
target name.

\begin{important}
Don't use \gls{gls} in chapter or section headings as it can have
some unpleasant side-effects. Instead use \gls{glsentrytext} for
regular entries and either \gls{glsentryshort} or \gls{glsentrylong}
for acronyms. Alternatively use
\sty{glossaries-extra} which provides special commands for use in
section headings and captions, such as \gls{glsfmttext} and \gls{glsfmtshort}.
\end{important}

The above examples are reasonably straightforward. The difficulty
comes if you want to display a \emph{sorted} list of all the entries you
have used in the document. The \sty{glossaries-extra} package
provides an easy way of listing all the defined entries:
\begin{codebox}
\cmd{documentclass}\marg{article}
\codepar
\cmd{usepackage}[\optval{sort}{none}]\marg{glossaries-extra}
\codepar
\gls{newglossaryentry}\marg{potato}\marg{\gloskeyval{name}{potato},
 \gloskeyval{plural}{potatoes},
 \gloskeyval{description}{starchy tuber}}
\codepar
\gls{newglossaryentry}\marg{cabbage}\marg{\gloskeyval{name}{cabbage},
 \gloskeyval{description}{vegetable with thick green
   or purple leaves}}
\codepar
\gls{newglossaryentry}\marg{turnip}\marg{\gloskeyval{name}{turnip},
  \gloskeyval{description}{round pale root vegetable}}
\codepar
\gls{newglossaryentry}\marg{carrot}\marg{\gloskeyval{name}{carrot},
  \gloskeyval{description}{orange root}}
\codepar
\cbeg{document}
Chop the \gls{gls}\marg{cabbage}, \gls{glspl}\marg{potato} and 
\gls{glspl}\marg{carrot}.
\codepar
\gls+{printunsrtglossaries} \comment{list all entries}
\cend{document}
\end{codebox}
However this method doesn't sort the entries (they're listed in
order of definition) and it will display all the defined entries,
regardless of whether or not you've used them all in the document,
so \qt{turnip} appears in the glossary even though there's
no \code{\gls{gls}\marg{turnip}} (or similar) in the document.

The \optval{sort}{none} option isn't essential in this case (there's
no other sort option available for this document), but it prevents
the automatic construction of the sort value and so slightly
improves the document build time.

This example document uses the same command
(\gls{printunsrtglossaries}) that's used with \app{bib2gls} (\option{b2g})
but with \app{bib2gls} you instead need to use the 
\opt{record} package option (not \opt{sort}) and one or more instances of
\gls{GlsXtrLoadResources} in the preamble (more on this later).

Most users prefer to have an automatically sorted list that only
contains entries that have been used in the document, optionally
with a page list (\gls{indexing}).
The \sty{glossaries} package provides
three options: use \TeX\ to perform the sorting (\option{noidx}); use
\app{makeindex} to perform the sorting (\option{mkidx}); use \app{xindy} to
perform the sorting (\option{xdy}). The extension package \sty{glossaries-extra}
provides a fourth method: use \app{bib2gls} (\option{b2g}).

The first option (using \TeX) is the simplest method, as it doesn't
require an external tool. It works best with at least version 3.0 of
the \sty{datatool-base} package, which is provided with
\sty{datatool} (and automatically loaded by \sty{glossaries}). 
That version introduced the ability to provide localisation support,
which needs to be installed separately. (See the \sty{datatool}
manual for further details.)

To use this method, add
\gls{makenoidxglossaries} to the preamble (before the entries are
defined) and put \gls{printnoidxglossaries} at the place where 
you want your glossary.
For example:
\begin{codebox}
\cmd{documentclass}\marg{article}
\codepar
\cmd{usepackage}\marg{glossaries}
\codepar
\gls{makenoidxglossaries} \comment{use TeX to sort}
\codepar
\gls{newglossaryentry}\marg{potato}\marg{\gloskeyval{name}{potato},
 \gloskeyval{plural}{potatoes},
 \gloskeyval{description}{starchy tuber}}
\codepar
\gls{newglossaryentry}\marg{cabbage}\marg{\gloskeyval{name}{cabbage},
 \gloskeyval{description}{vegetable with thick green
   or purple leaves}}
\codepar
\gls{newglossaryentry}\marg{turnip}\marg{\gloskeyval{name}{turnip},
 \gloskeyval{description}{round pale root vegetable}}
\codepar
\gls{newglossaryentry}\marg{carrot}\marg{\gloskeyval{name}{carrot},
 \gloskeyval{description}{orange root}}
\codepar
\cbeg{document}
Chop the \gls{gls}\marg{cabbage}, \gls{glspl}\marg{potato} and 
\gls{glspl}\marg{carrot}.
\codepar
\gls{printnoidxglossaries}
\cend{document}
\end{codebox}

Try this out and run \LaTeX\ (or \pdfLaTeX) \emph{twice}. The first
run won't show the \idx{glossary}. It will only appear on the second run.
This doesn't include \qt{turnip} in the \idx{glossary} because that term
hasn't been referenced (with commands like \code{\gls{gls}\marg{turnip}}) 
in the document.

The glossary has a vertical gap between the \qt{carrot} term and the
\qt{potato} term. This is because the entries in the glossaries are
grouped according to their first letter. If you don't want this gap,
just add \opt{nogroupskip} to the package options:
\begin{codebox}
\cmd{usepackage}[\opt{nogroupskip}]\marg{glossaries}
\end{codebox}
or you may want to try out a style that shows the group headings:
\begin{codebox*}
\cmd{usepackage}[\optval{style}{\glostyle{indexgroup}}]\marg{glossaries}
\end{codebox*}
If you try out this example you may also notice that the description
is followed by a full stop (period) and a number. The number is the
location in the document where the entry was referenced (page~1 in this
case), so you can look up the term in the \idx{glossary} and be directed to
the relevant pages. It may be that you don't want this
back-reference, in which case you can suppress it using the
\opt{nonumberlist} package option:
\begin{codebox}
\cmd{usepackage}[\opt{nonumberlist}]\marg{glossaries}
\end{codebox}
If you don't like the terminating full stop, you can suppress that
with the \opt{nopostdot} package option:
\begin{codebox}
\cmd{usepackage}[\opt{nopostdot}]\marg{glossaries}
\end{codebox}
The \sty{glossaries-extra} package has this option as the default,
so if you want terminating full stop, you need to explicitly add it
with the \opt{postdot} option:
\begin{codebox}
\cmd{usepackage}[\opt{postdot}]\marg{glossaries-extra}
\end{codebox}

You may have noticed that I've used another command in the above examples:
\glsadd{glslike}\gls{glspl}. This displays the plural form. By default, this is just
the singular form with the letter \qt{s} appended, but in the case of
\qt{potato} I had to specify the correct plural using the
\gloskey{plural} key.

Using \TeX\ to sort the entries is the simplest but least efficient
method. If you have a large \idx{glossary}, then you will have a
much faster build time if you use \app{makeindex} (\option{mkidx})
or \app{xindy} (\option{xdy}) or \app{bib2gls} (\option{b2g}). If
you are using extended Latin or non-Latin characters, then
\app{xindy} or \app{bib2gls} are the recommended methods.  These
methods are described in more detail in
\sectionref{sec:printglossaries}.

The rest of this document briefly describes the main commands
provided by the \sty{glossaries} package. (Most of these are also
available with \sty{glossaries-extra} but may behave slightly
differently.)

\section{Defining Terms}
\label{sec:defterm}

When you use the \sty{glossaries} package, you need to define
glossary entries before you can reference them. This is best 
done in the document preamble, as shown in the earlier examples, or
in a separate file that's input in the preamble.

These entries could be a~word, phrase, abbreviation or symbol.
They're usually accompanied by a~description, which could be a short
sentence or an in-depth explanation that spans multiple paragraphs.
The simplest method of defining an entry is to use:
\cmddef{newglossaryentry}
where \meta{label} is a unique label that identifies this entry.
(Don't include the angle brackets \meta{ }. They just indicate the parts of
the code that you need to change when you use this command in your document.) 
The two most important keys are \gloskey{name} and
\gloskey{description}:
\begin{compactcodebox}
\gls{newglossaryentry}\margm{label}
\marg{
  \gloskeyval{name}{\meta{name}},
  \gloskeyval{description}{\meta{description}},
  \meta{other options}
}
\end{compactcodebox}%
The \meta{name} is the word, phrase or symbol you are defining,
and \meta{description} is the description to be displayed in the
glossary.

This command is a \qt{short} command, which means that
\meta{description} can't contain a~paragraph break. If you have
a~long description, you can instead use:
\cmddef{longnewglossaryentry}
In this case the \gloskey{name} key is in the second argument but the
description is supplied in the third argument instead of via a key.
\begin{compactcodebox}
\gls{longnewglossaryentry}\margm{label}
\marg{
  \gloskeyval{name}{\meta{name}},
  \meta{other options}
}
\margm{description}
\end{compactcodebox}

Examples:
\begin{enumerate}
\item Define the term \qt{set} with the label \code{set}:
\begin{codebox}
\gls{newglossaryentry}\marg{set}
\marg{
  \gloskeyval{name}{set},
  \gloskeyval{description}{a collection of objects}
}
\end{codebox}

\item Define the symbol $\emptyset$ with the label
\code{emptyset}:
\begin{codebox}
\gls{newglossaryentry}\marg{emptyset}
\marg{
  \gloskeyval{name}{\cmd{ensuremath}\marg{\cmd{emptyset}}},
  \gloskeyval{description}{the empty set}
}
\end{codebox}
(This will also need a \gloskey{sort} key if you use \optionsor{noidx,xdy}, see
below.)

\item Define the phrase \qt{Fish Age} with the label
\code{fishage}:
\begin{codebox}
\gls{longnewglossaryentry}\marg{fishage}
\marg{\gloskeyval{name}{Fish Age}}
\marg{\%
  A common name for the Devonian geologic period 
  spanning from the end of the Silurian Period to
  the beginning of the Carboniferous Period.
\codepar
  This age was known for its remarkable variety of 
  fish species.
}
\end{codebox}
(The percent character discards the end of line character that would
otherwise cause an unwanted space to appear at the start of the
description.)

\item If you are using UTF-8 characters with the \sty{inputenc}
package, make sure you have \sty{mfirstuc} v2.08+ installed:
\begin{codebox}
\comment{\sty{mfirstuc} v2.08+}
\gls{newglossaryentry}\marg{\'elite}
\marg{
  \gloskeyval{name}{\'elite},
  \gloskeyval{description}{select group or class}
}
\end{codebox}
If you have an older version of \sty{mfirstuc} then any initial
character that is an extended Latin or non-Latin character must be
grouped in order to work with sentence-casing commands, such as
\gls{Gls}. An older version possibly also means there may not be support
for UTF-8 characters in labels either:
\begin{codebox}
\comment{\sty{mfirstuc} v2.07 or older}
\gls{newglossaryentry}\marg{elite}
\marg{
  \gloskeyval{name}{\marg{\'e}lite},
  \gloskeyval{description}{select group or class}
}
\end{codebox}
Likewise if accent commands are used:
\begin{codebox}
\gls{newglossaryentry}\marg{elite}
\marg{
  \gloskeyval{name}{\marg{\cmd{'}e}lite},
  \gloskeyval{description}{select group or class}
}
\end{codebox}
\end{enumerate}

If you use \app{bib2gls} with \sty{glossaries-extra} then the
terms must be defined in a \ext+{bib} file. For example:
\begin{codebox}
\comment{Encoding: UTF-8}
\codepar
\atentry{entry}\marg{set,
  \gloskeyval{name}{set},
  \gloskeyval{description}{a collection of objects}
}
\codepar
\atentry{entry}\marg{emptyset,
  \gloskeyval{name}{\cmd{ensuremath}\marg{\cmd{emptyset}}},
  \gloskeyval{description}{the empty set}
}
\codepar
\atentry{entry}\marg{fishage,
  \gloskeyval{name}{Fish Age},
  \gloskeyval{description}{A common name for the Devonian 
  geologic period spanning from the end of the 
  Silurian Period to the beginning of the 
  Carboniferous Period.
\codepar
  This age was known for its remarkable variety of 
  fish species.}
}
\codepar
\atentry{entry}\marg{\'elite,
  \gloskeyval{name}{\'elite},
  \gloskeyval{description}{select group or class}
}
\end{codebox}
(The \ext+{bib} format doesn't allow spaces in labels so you can't
have \code{fish age} as the label, but you can have 
\code{fish-age}.)
This method requires \sty{glossaries-extra}['s] \opt{record}
package option:
\begin{codebox}
\cmd{usepackage}[\opt{record}]\marg{glossaries-extra}
\end{codebox}
The \ext{bib} file is specified in the resource command. For
example, if the \ext{bib} file is called \filefmt{entries.bib}
then put the following line in the document preamble:
\begin{codebox}
\gls{GlsXtrLoadResources}\oarg{\resourceoptvalm{src}{entries}}
\end{codebox}
You can have a comma-separated list. For example, if you also
have entries defined in the file \filefmt{entries2.bib}:
\begin{codebox}
\gls{GlsXtrLoadResources}\oarg{\resourceoptvalm{src}{entries,entries2}}
\end{codebox}

There are other keys you can use when you define an entry. For
example, the \gloskey{name} key indicates how the term
should appear in the list of entries (\idx{glossary}), but if the term should
appear differently when you reference it with \code{\gls{gls}\margm{label}} in 
the document, you need to use the \gloskey{text} key as well.

For example:
\begin{codebox}
\gls{newglossaryentry}\marg{latinalph}
\marg{
  \gloskeyval{name}{alphabet, Latin},
  \gloskeyval{text}{Latin alphabet},
  \gloskeyval{description}{alphabet consisting of the letters 
  a, \cmd{ldots}, z, A, \cmd{ldots}, Z}
}
\end{codebox}
This will appear in the text as \qt{Latin alphabet} but will be listed in
the glossary as \qt{alphabet, Latin} (and therefore sorted under
\qt{A} not \qt{L}).
With \app{bib2gls} this entry is defined in the \ext{bib}
file as:
\begin{codebox}
\atentry{entry}\marg{latinalph,
  \gloskeyval{name}{alphabet, Latin},
  \gloskeyval{text}{Latin alphabet},
  \gloskeyval{description}{alphabet consisting of the letters 
  a, \cmd{ldots}, z, A, \cmd{ldots}, Z}
}
\end{codebox}

Another commonly used key is \gloskey{plural} for specifying the
plural of the term. This defaults to the value of the \gloskey{text}
key with an \qt{s} appended, but if this is incorrect, just use the
\gloskey{plural} key to override it:
\begin{codebox}
\gls{newglossaryentry}\marg{oesophagus}
\marg{
  \gloskeyval{name}{\oe{}sophagus},
  \gloskeyval{plural}{\oe{}sophagi},
  \gloskeyval{description}{canal from mouth to stomach}
}
\end{codebox}

Abbreviations can be defined using:
\cmddef{newacronym}
where \meta{label} is the label (as per \gls{newglossaryentry}), 
\meta{short} is the short form and \meta{long} is the long form. 
For example, the
following defines an abbreviation:
\begin{codebox}
\gls{newacronym}\marg{svm}\marg{SVM}\marg{support vector machine}
\end{codebox}
This internally uses \gls{newglossaryentry} to define an entry with
the label \code{svm}. By default, the \gloskey{name} key is set to
\meta{short} (\qt{SVM} in the above example) and the
\gloskey{description} key is set to \meta{long} (\qt{support vector
machine} in the above example). If, instead, you want to be able to
specify your own description you can do this using the optional
argument:
\begin{codebox}
\gls{newacronym}
 \oarg{\gloskeyval{description}{statistical pattern recognition
 technique}}
 \marg{svm}\marg{SVM}\marg{support vector machine}
\end{codebox}

Before you define your acronyms, you need to specify which style to use with:
\cmddef{setacronymstyle}
where \meta{style-name} is the name of the style. There are a number
of predefined styles, such as: \acrstyle{long-short} (on \idx{firstuse}
display the long form with the short form in parentheses);
\acrstyle{short-long} (on \idx{firstuse} display the short form with the
long form in parentheses); \acrstyle{long-short-desc} (like
\acrstyle{long-short} but you need to specify the description); or
\acrstyle{short-long-desc} (like \acrstyle{short-long} but you need to
specify the description). There are some other styles as well that
use \gls{textsc} to typeset the acronym in small-caps or that use a footnote on
first use. See the main user guide for further details.

The \sty{glossaries-extra} package provides improved abbreviation
handling with \gallerypage{sample-abbr-styles}{a lot more predefined styles}.
With this extension package, abbreviations are defined using:
\cmddef{newabbreviation}
You can still use \gls{newacronym} but it's redefined to 
use the new abbreviation interface. The style must now be set using:
\cmddef{setabbreviationstyle}
The default \meta{category} is \cat{abbreviation}. If you use
\gls{newacronym} the category is \cat{acronym}, which is why you
need to use the optional argument if you define abbreviations with
\gls{newacronym} when \sty{glossaries-extra} has been loaded:
\begin{compactcodebox*}
\gls{setabbreviationstyle}\oarg{\cat{acronym}}\margm{style name}
\end{compactcodebox*}
If you use \app{bib2gls} then abbreviations are defined in the
\ext{bib} file in the format:
\begin{compactcodebox*}
\atentry{abbreviation}\marg{\meta{label},
  \gloskeyval{long}{\meta{long form}},
  \gloskeyval{short}{\meta{short form}}
}
\end{compactcodebox*}

The plural forms for abbreviations can be specified using the
\gloskey{longplural} and \gloskey{shortplural} keys. For example:
\begin{codebox}
\gls{newacronym}
 \oarg{\gloskeyval{longplural}{diagonal matrices}}
 \marg{dm}\marg{DM}\marg{diagonal matrix}
\end{codebox}
or (with \sty{glossaries-extra}):
\begin{codebox}
\gls{newabbreviation} \comxr
 \oarg{\gloskeyval{longplural}{diagonal matrices}}
 \marg{dm}\marg{DM}\marg{diagonal matrix}
\end{codebox}
If omitted, the defaults are again obtained by appending an \qt{s} to
the singular versions. (The \sty{glossaries-extra} package provides
a way of not appending \qt{s} for abbreviation plurals via category
attributes.) With \app{bib2gls}, the definition in the
\ext{bib} file is:
\begin{codebox}
\atentry{abbreviation}\marg{dm,
  \gloskeyval{short}{DM},
  \gloskeyval{long}{diagonal matrix},
  \gloskeyval{longplural}{diagonal matrices}
}
\end{codebox}

It's also possible to have both a~name and a~corresponding symbol.
Just use the \gloskey{name} key for the name and the \gloskey{symbol} key
for the symbol. For example:
\begin{codebox}
\gls{newglossaryentry}\marg{emptyset}
\marg{
  \gloskeyval{name}{empty set},
  \gloskeyval{symbol}{\cmd{ensuremath}\marg{\cmd{emptyset}}},
  \gloskeyval{description}{the set containing no elements}
}
\end{codebox}
or with \app{bib2gls} the definition in the \ext{bib} file is:
\begin{codebox}
\atentry{entry}\marg{emptyset,
  \gloskeyval{name}{empty set},
  \gloskeyval{symbol}{\cmd{ensuremath}\marg{\cmd{emptyset}}},
  \gloskeyval{description}{the set containing no elements}
}
\end{codebox}
If you want commands such as \csfmt{emptyset} in the \gloskey{name} field 
then you must supply a \gloskey{sort} value with \options{noidx,xdy} 
otherwise you'll end up with errors from \TeX\ or \app{xindy}. With \option{mkidx}
(\app{makeindex}) you're unlikely to get an error, but you may find
the resulting order is a little odd. For example:
\begin{codebox}
\gls{newglossaryentry}\marg{emptyset}
\marg{
  \gloskeyval{name}{\cmd{ensuremath}\marg{\cmd{emptyset}}},
  \gloskeyval{sort}{empty set},
  \gloskeyval{description}{the set containing no elements}
}
\end{codebox}
This displays the symbol as $\emptyset$ but sorts according to
\qt{empty set}. You may want to consider using
\sty{glossaries-extra}['s] \opt{symbols} package option which
provides
\cmddef{glsxtrnewsymbol}
This internally uses \gls{newglossaryentry} but automatically sets
the \gloskey{sort} key to the \meta{label}. For example:
\begin{codebox}
\cmd{documentclass}\marg{article}
\cmd{usepackage}[\opt{symbols}]\marg{glossaries-extra}
\gls{makeglossaries}
\gls{glsxtrnewsymbol}
 \oarg{\gloskeyval{description}{the set containing no elements}}
 \marg{emptyset}\comment{label (and sort value)}
 \marg{\cmd{ensuremath}\marg{\cmd{emptyset}}}\comment{symbol}
\cbeg{document}
\gls{gls}\marg{emptyset}
\gls{printglossaries}
\cend{document}
\end{codebox}
Now the sort value is automatically set to the label \qt{emptyset}.

With \app{bib2gls} you can define this entry in the \ext{bib} file
as
\begin{codebox}
\atentry{entry}\marg{emptyset,
  \gloskeyval{name}{\cmd{ensuremath}\marg{\cmd{emptyset}}},
  \gloskeyval{description}{the set containing no elements}
}
\end{codebox}
in which case \app{bib2gls} will try to interpret the \gloskey{name} 
field to determine the sort value. Alternatively you can use:
\begin{codebox}
\strong{\atentry{symbol}}\marg{emptyset,
  \gloskeyval{name}{\cmd{ensuremath}\marg{\cmd{emptyset}}},
  \gloskeyval{description}{the set containing no elements}
}
\end{codebox}
which will use the label (\code{emptyset}) as the sort value.
(You don't need the \opt{symbols} package option in this case,
unless you want a separate symbols list.) The corresponding document
(where the definition is in the file \filefmt{entries.bib}) is now:
\begin{codebox}
\cmd{documentclass}\marg{article}
\cmd{usepackage}[\opt{record}]\marg{glossaries-extra}
\gls{GlsXtrLoadResources}\oarg{\resourceoptvalm{src}{entries}}
\cbeg{document}
\gls{gls}\marg{emptyset}
\gls{printunsrtglossaries}
\cend{document}
\end{codebox}

Note that, while the \gloskey{sort} key is advised for symbols when
using \gls{makeglossaries} or \gls{makenoidxglossaries}, the
\gloskey{sort} field shouldn't be used with \app{bib2gls}. Instead,
\app{bib2gls} has its own algorithm for determining the sort value
based on the entry type (\atentry{entry}, \atentry{symbol} etc) and
various settings. See
\gallerypage{bib2gls-sorting}{bib2gls gallery: sorting} for further
details.

\section{Using Entries}
\label{sec:useterm}

Once you have defined your entries, as described above, you can
reference them in your document. There are a~number of commands to
do this, but the most common one is:
\cmddef{gls}
where \meta{label} is the label you assigned to the entry when you
defined it. For example, \code{\gls{gls}\marg{fishage}} will display \qt{Fish
Age} in the text (given the definition from the previous section).
If you are using \app{bib2gls} then this will display \idx{unknowntag} (like
\csfmt{ref} and \csfmt{cite}) until \app{bib2gls} has created the
relevant files and \LaTeX\ is rerun.

If you are using the \sty{hyperref} package (remember to load it
before \sty{glossaries}) \gls{gls} will create a hyperlink to the
corresponding entry in the glossary. If you want to suppress the
hyperlink for a particular instance, use the starred form \code{\gls{gls}*}
for example, \code{\gls{gls}*\marg{fishage}}. The other commands described in
this section all have a similar starred form.

If the entry was defined as an acronym (using \gls{newacronym} with
\sty{glossaries} described above) or an abbreviation (using
\gls{newabbreviation} with \sty{glossaries-extra}), then \gls{gls}
will display the full form the first time it's used (\idx{firstuse}) and just the
short form on \idx{subsequentuse}. For example, if the style is
set to \acrstyle{long-short}, then \code{\gls{gls}\marg{svm}} will display
\qt{support vector machine (SVM)} the first time it's used, but the
next occurrence of \code{\gls{gls}\marg{svm}} will just display \qt{SVM}.
(If you use \gls{newacronym} with \sty{glossaries-extra} the default
doesn't show the long form on \idx{firstuse}. You'll need to change the
style first, as described earlier.)

If you want the plural form, you can use:
\cmddef{glspl}
instead of \code{\gls{gls}\margm{label}}. For example,
\code{\gls{glspl}\marg{set}} displays \qt{sets}.

If the term appears at the start of a~sentence, you can convert the
first letter to \idx{uppercase} (\idx{sentencecase}) using:
\cmddef{Gls}
for the singular form or
\cmddef{Glspl}
for the plural form. For example:
\begin{codebox}
\gls{Glspl}\marg{set} are collections.
\end{codebox}
produces \qt{Sets are collections}.

If you've specified a symbol using the \gloskey{symbol} key, you can
display it using:
\cmddef{glssymbol}
For example
\begin{codebox}
\cmd{documentclass}\marg{article}
\cmd{usepackage}\marg{glossaries}
\gls{newglossaryentry}\marg{emptyset}
\marg{
  \gloskeyval{name}{empty set},
  \gloskeyval{symbol}{\cmd{ensuremath}\marg{\cmd{emptyset}}},
  \gloskeyval{description}{the set containing no elements}
}
\cbeg{document}
The \gls{gls}\marg{emptyset} is denoted \gls{glssymbol}\marg{emptyset}.
\cend{document}
\end{codebox}
\begin{resultbox}
The empty set is denoted $\emptyset$.
\end{resultbox}

\section{Displaying a List of Entries}
\label{sec:printglossaries}

In \sectionref{sec:start}, I mentioned that there are three options
you can choose from to create an automatically sorted glossary with the base 
\sty{glossaries} package. These are also available with the
extension package \sty{glossaries-extra} along with a fourth
option. These four options are listed below in a little more detail.
\Tableref{tab:optionsp+c} summarises the main advantages and
disadvantages.  (There's a more detailed summary in the main
\sty{glossaries} user manual.) See also
\dickimawhref{latex/buildglossaries/}{Incorporating makeglossaries or makeglossaries-lite or bib2gls into the document build}.

Note that \app{makeindex} and \app{xindy} are general purpose
indexing applications that are developed independently of
\sty{glossaries} and \sty{glossaries-extra}. This makes them harder
to integrate, as glossaries have more complex requirements than a
simple index, particularly for documents where some terms may only
appear in the description of another term. (There's also the
possibility that they may change in a manner that becomes
incompatible with \sty{glossaries} or may even be declared obsolete.)
The \app{bib2gls} application, on the other hand, is designed
specifically for, and developed alongside, the
\sty{glossaries-extra} package.

\begin{table}[htbp]
 \caption{Comparison of Glossary Options}
 \label{tab:optionsp+c}
 {%
 \centering
 \begin{tabular}{>{\raggedright}p{0.3\textwidth}cccc}
   & \bfseries \option{noidx} 
   & \bfseries \option{mkidx} 
   & \bfseries \option{xdy}
   & \bfseries \option{b2g}\\
   Requires \sty{glossaries-extra}? &
   \conno & \conno & \conno & \conyes\\
   Requires an external application? &
   \conno & \conyes & \conyes & \conyes\\
   Requires Perl? &
   \conno & \conno & \conyes & \conno\\
   Requires Java? &
   \conno & \conno & \conno & \conyes\\
   Can sort extended Latin
   or non-Latin alphabets? &
   \prono\tablefnmark{\asteriskmarker} & \prono & \proyes & \proyes\\
   Efficient sort algorithm? &
   \prono & \proyes & \proyes & \proyes\\
   Can use different sort methods for each glossary? &
   \proyes & \prono & \prono & \proyes\\
   Any problematic sort values? &
   \conyes & \conyes & \conyes & \conno\\
   Can form ranges in the location lists? &
   \prono & \proyes & \proyes & \proyes\\
   Can have non-standard locations? &
   \proyes & \prono & \proyes\tablefnmark{\textdagger} & \proyes
 \end{tabular}
 \par}\medskip
 \tablefntext{\asteriskmarker}{Requires corresponding \sty{datatool} v3.0+
localisation support.}
 \tablefntext{\textdagger}{Requires some setting up.}
\end{table}

\begin{deflist}
\itemtitle{\idxoptiondef{noidx}}

\begin{itemdesc}
 This is the simplest option but it's slow and if you want a sorted list,
 it can be problematic with extended or non-Latin alphabets.
 If the \gloskey{name} contains commands then  the
 \gloskey{sort} value must be supplied unless you 
 using the option \optval{sort}{def} 
 or \optval{sort}{use}. This method works best with at least version
 3.0 of the \sty{datatool-base} package (which is supplied with
 \sty{datatool} and automatically loaded by \sty{glossaries}).

  \begin{enumerate}
    \item Add \gls{makenoidxglossaries} to your preamble (before you
    start defining your entries, as described in
    \sectionref{sec:defterm}).

    \item Put
\begin{compactcodebox*}
\gls{printnoidxglossary}\oarg{\printglossoptval{sort}{\meta{order}},\meta{other options}}
\end{compactcodebox*}
    where you want your list of entries to appear. The sort
    \meta{order} may be one of: \code{word} (word ordering),
    \code{letter} (letter ordering), \code{case} (case-sensitive
    ordering), \code{def} (in order of definition) or
    \code{use} (in order of use). Alternatively, use
\begin{compactcodebox*}
\gls{printnoidxglossaries}
\end{compactcodebox*}%
    to display all your glossaries (if you have more than one).

    This option allows you to have different sort methods. For
    example:
\begin{codebox}
\gls{printnoidxglossary}\oarg{\printglossoptval{sort}{word}}\comment{main glossary}
\gls{printnoidxglossary}\oarg{
 \printglossoptval{type}{symbols},\comment{symbols glossary}
 \printglossoptval{sort}{use}}
\end{codebox}

    \item Run \LaTeX\ twice on your document. (As you would do to
    make a~table of contents appear.) For example, click twice on
    the \qt{typeset} or \qt{build} or \qt{PDF\LaTeX} button in your editor.
  \end{enumerate}
Here's a complete document (\filefmt{myDoc.tex}):
\begin{codebox}
\cmd{documentclass}\marg{article}
\codepar
\cmd{usepackage}\marg{glossaries}
\codepar
\gls{makenoidxglossaries}\comment{use TeX to sort}
\codepar
\gls{newglossaryentry}\marg{sample}\marg{\gloskeyval{name}{sample},
  \gloskeyval{description}{an example}}
\codepar
\cbeg{document}
A \gls{gls}\marg{sample}.
\codepar
\comment{iterate over all indexed entries:}
\gls{printnoidxglossaries}
\cend{document}
\end{codebox}
Document build:
\begin{terminal}
pdflatex myDoc
pdflatex myDoc
\end{terminal}
\end{itemdesc}

\itemtitle{\idxoptiondef{mkidx}}

\begin{itemdesc}
   This option uses an application called \app{makeindex} to sort 
   the entries. This application comes with all modern \TeX\ distributions, 
   but it's hard-coded for the non-extended Latin alphabet. This process 
   involves making \LaTeX\ write the glossary information to a temporary 
   file which \app{makeindex} reads. Then \app{makeindex} writes 
   a~new file containing the code to typeset the glossary. \LaTeX\ then 
   reads this file on the next run. The \app{makeindex}
   application is automatically invoked by the helper
   \app{makeglossaries} script, which works out all the
   appropriate settings from the \ext+{aux} file.

   \begin{enumerate}
    \item Add \gls{makeglossaries} to your preamble (before you start
    defining your entries).

    \item Put
\begin{compactcodebox*}
\gls{printglossary}\oargm{options}
\end{compactcodebox*}%
    where you want your list of entries (\idx{glossary}) to appear. (The
    \printglossopt{sort} key isn't available in \meta{options}.)
    Alternatively, use
\begin{compactcodebox*}
\gls{printglossaries}
\end{compactcodebox*}%
    which will display all glossaries (if you have more than one).

\begin{important}
All glossaries are sorted using the same method
which may be identified with one of the package options:
\optval{sort}{standard} (default), \optval{sort}{use}
or \optval{sort}{def}.
\end{important}

    \item Run \LaTeX\ on your document. This creates files with the
    extensions \ext+{glo} and \ext+{ist} (for example, if your 
    \LaTeX\ document is called \filefmt{myDoc.tex}, then you'll have 
    two extra files called \filefmt{myDoc.glo} and \filefmt{myDoc.ist}).
    If you look at your document at this point, you won't see the 
    \idx{glossary} as it hasn't been created yet.

    \item Run \app{makeglossaries} with the base name of your 
    document (without the \ext{tex} extension). If
    you have access to a terminal or a command prompt then you need 
    to run the command:
\begin{terminal}
makeglossaries myDoc
\end{terminal}%
   (Replace \filefmt{myDoc} with the base name of your \LaTeX\
    document file without the extension.) If you don't have Perl installed
    use \app{makeglossaries-lite} instead:
\begin{terminal}
makeglossaries-lite myDoc
\end{terminal}

\begin{important}
Some beginners get confused by \app{makeglossaries} the
application (run as a system command) and \gls{makeglossaries}
the \LaTeX\ command which should be typed in the document
preamble. These are two different concepts that happen to have
similar looking names.
\end{important}

    If you don't know how to use the command prompt, then you can 
    probably configure your text editor to add 
    \app{makeglossaries} (or \app{makeglossaries-lite})
    as a build tool, but each editor has a
    different method of doing this, so I~can't give a~general
    description. However, there are some guidelines in
\dickimawhref{latex/buildglossaries/}{Incorporating
makeglossaries or makeglossaries-lite or bib2gls into the document
build}.
    If you still have problems, try adding the \opt{automake}
    package option:
\begin{compactcodebox}
\cmd{usepackage}[\opt{automake}]\marg{glossaries}
\end{compactcodebox}

    The default sort is word order (\qt{sea lion} comes before \qt{seal}). 
    If you want letter ordering you need to add the
    \optval{order}{letter} package option
\begin{compactcodebox}
\cmd{usepackage}[\optval{order}{letter}]\marg{glossaries}
\end{compactcodebox}

    \item Once you have successfully completed the previous step,
    you can now run \LaTeX\ on your document again.
   \end{enumerate}
Here's a complete document (\filefmt{myDoc.tex}):
\begin{codebox}
\cmd{documentclass}\marg{article}
\codepar
\cmd{usepackage}\marg{glossaries}
\codepar
\gls{makeglossaries} \comment{create \app{makeindex} files}
\codepar
\gls{newglossaryentry}\marg{sample}\marg{\gloskeyval{name}{sample},
  \gloskeyval{description}{an example}}
\codepar
\cbeg{document}
A \gls{gls}\marg{sample}.
\codepar
\comment{input files created by makeindex:}
\gls{printglossaries}
\cend{document}
\end{codebox}
Document build:
\begin{terminal}
pdflatex myDoc
makeglossaries myDoc
pdflatex myDoc
\end{terminal}
or
\begin{terminal}
pdflatex myDoc
makeglossaries-lite myDoc
pdflatex myDoc
\end{terminal}
\end{itemdesc}

\itemtitle{\idxoptiondef{xdy}}

\begin{itemdesc}
   This option uses an application called
   \app{xindy} to sort the entries. This application is more
   flexible than \app{makeindex} and is able to sort extended
   Latin or non-Latin alphabets. It comes with both \TeX~Live and 
   MiK\TeX. Since \app{xindy} is a Perl script, you will also 
   need to ensure that Perl is installed. In a~similar way to \option{mkidx}, this 
   option involves making \LaTeX\ write the glossary information to 
   a~temporary file which \app{xindy} reads. Then \app{xindy} 
   writes a~new file containing the code to typeset the glossary. 
   \LaTeX\ then reads this file on the next run. The \app{xindy}
   application is automatically invoked by the helper
   \app{makeglossaries} script, which works out all the
   appropriate settings from the \ext+{aux} file.
   See \url{http://www.xindy.org/} for further information about
   \app{xindy} and links to its mailing list and issue tracker.

   \begin{enumerate}
     \item Add the \opt{xindy} option to the \sty{glossaries}
package option list:
\begin{compactcodebox}
\cmd{usepackage}[\opt{xindy}]\marg{glossaries}
\end{compactcodebox}%
If you aren't using a Latin script, you may need to suppress the
default number group:
\begin{compactcodebox}
\cmd{usepackage}[\optvalm{xindy}{glsnumbers=false}]\marg{glossaries}
\end{compactcodebox}

     \item Add \gls{makeglossaries} to your preamble (before you start
     defining your entries).

    \item Put
\begin{compactcodebox*}
\gls{printglossary}\oargm{options}
\end{compactcodebox*}%
    where you want your list of entries (glossary) to appear. (The
    \printglossopt{sort} key isn't available in \meta{options}.)
    Alternatively, use
\begin{compactcodebox*}
\gls{printglossaries}
\end{compactcodebox*}

\begin{important}
All \idxpl{glossary} are sorted using the same method
which may be identified with one of the package options:
\optval{sort}{standard} (default), \optval{sort}{use}
or \optval{sort}{def}.
\end{important}

    \item Run \LaTeX\ on your document. This creates files with the
    extensions \ext+{glo} and \ext+{xdy} (for example, if your 
    \LaTeX\ document is called \filefmt{myDoc.tex}, then you'll have 
    two extra files called \filefmt{myDoc.glo} and \filefmt{myDoc.xdy}).
    If you look at your document at this point, you won't see the 
    \idx{glossary} as it hasn't been created yet.

    \item Run \app{makeglossaries} with the base name of the
    document (omitting the \ext{tex} extension).
    If you have access to a terminal or a command prompt then 
    you need to run the command:
\begin{terminal}
makeglossaries myDoc
\end{terminal}%
    (Replace \filefmt{myDoc} with the base name of your \LaTeX\
    document file without the \ext{tex} extension.)
    If you don't know how to use the command prompt, then as mentioned above, you may
    be able to configure your text editor to add
    \app{makeglossaries} as a build tool. Note that the \opt{automake} option
    won't work in \TeX's restricted mode, as \app{xindy} isn't on the
    list of trusted applications.

    The default sort is word order (\qt{sea lion} comes before \qt{seal}). 
    If you want letter ordering you need to add the
    \optval{order}{letter} package option:
\begin{compactcodebox}
\cmd{usepackage}[\opt{xindy},\optval{order}{letter}]\marg{glossaries}
\end{compactcodebox}

    \item Once you have successfully completed the previous step,
    you can now run \LaTeX\ on your document again.

   \end{enumerate}
Here's a complete document (\filefmt{myDoc.tex}):
\begin{codebox}
\cmd{documentclass}\marg{article}
\codepar
\cmd{usepackage}[\opt{xindy}]\marg{glossaries}
\codepar
\gls{makeglossaries} \comment{create \app{xindy} files}
\codepar
\gls{newglossaryentry}\marg{sample}\marg{\gloskeyval{name}{sample},
  \gloskeyval{description}{an example}}
\codepar
\cbeg{document}
A \gls{gls}\marg{sample}.
\codepar
\gls{printglossaries} \comment{input files created by \app{xindy}}
\cend{document}
\end{codebox}%
Document build:
\begin{terminal}
pdflatex myDoc
makeglossaries myDoc
pdflatex myDoc
\end{terminal}

If you use \app{xindy} and the \gloskey{name} only contains a command or
commands (such as \csfmt{P} or
\code{\cmd{ensuremath}\marg{\cmd{pi}}}) you must add the
\gloskey{sort} key. This is also advisable for the other options
(except \option{b2g}, \gallerypage{bib2gls-sorting}{which, in general, shouldn't
have the \optfmt{sort} field set}), but is essential for \option{xdy}. For example:
\begin{codebox}
\gls{newglossaryentry}\marg{P}\marg{\gloskeyval{name}{\cmd{P}},\gloskeyval{sort}{P},
  \gloskeyval{description}{paragraph symbol}}
\end{codebox}

\end{itemdesc}

\itemtitle{\idxoptiondef{b2g}}

\begin{itemdesc}
   This requires the extension package \sty+{glossaries-extra} and
   an application called \app{bib2gls}. This application is able
   to sort extended Latin or non-Latin alphabets. It comes with both
   \TeX~Live and MiK\TeX\ but requires Java. This method
   works differently to \options{mkidx,xdy}. Instead of creating a file
   containing the code to typeset the glossary it creates a
   \ext+{glstex} file containing the entry definitions obtained from
   data that was fetched from the \ext+{bib} file (or files), but only those entries 
   that are required in the document are defined and they are 
   defined in the order obtained from the chosen sort method. 
   This means that you can just use
   \gls{printunsrtglossary} to display each glossary (or
   \gls{printunsrtglossaries} to display them all).

   \begin{enumerate}
   \item Add the \opt{record} option to the \sty{glossaries-extra}
   package option list:
\begin{compactcodebox}
\cmd{usepackage}\oarg{\opt{record}}\marg{glossaries-extra}
\end{compactcodebox}

   \item Add one or more 
\begin{compactcodebox*}
\gls{GlsXtrLoadResources}\oarg{\resourceoptvalm{src}{\meta{bib list}},\meta{options}}
\end{compactcodebox*}%
   to your preamble where \meta{bib list} is the
   list of \ext{bib} files containing the entries. You may
   use different sort methods for each resource set. For example:
\begin{codebox*}
\cmd{usepackage}\oarg{\opt{record},\comment{using \app{bib2gls}}
 \opt{abbreviations},\comment{create abbreviations list}
 \opt{symbols},\comment{create symbols list}
 \opt{numbers}\comment{create numbers list}
}\marg{glossaries-extra}
\codepar
\gls{GlsXtrLoadResources}\oarg{
  \resourceoptvalm{src}{terms},\comment{entries in terms.bib}
\comment{put these entries in the 'main' list:}
  \resourceoptval{type}{main},
\comment{sort according to this locale:}
  \resourceoptval{sort}{de-CH-1996}
}
\gls{GlsXtrLoadResources}\oarg{
  \resourceoptvalm{src}{abbrvs},\comment{entries in abbrvs.bib}
\comment{put these entries in the 'abbreviations' list:}
  \resourceoptvalm{type}{abbreviations},
\comment{case-sensitive letter (non-locale) sort:}
  \resourceoptval{sort}{letter-case}
}
\gls{GlsXtrLoadResources}\oarg{
  \resourceoptvalm{src}{syms},\comment{entries in syms.bib}
\comment{put these entries in the 'symbols' list:}
  \resourceoptval{type}{symbols},
\comment{sort according to first use in the document:}
  \resourceoptval{sort}{use}
}
\gls{GlsXtrLoadResources}\oarg{
  \resourceoptvalm{src}{constants},\comment{entries in constants.bib}
\comment{put these entries in the 'numbers' list:}
  \resourceoptval{type}{numbers},
\comment{sort according to the \gloskey{user1} field:}
  \resourceoptval{sort-field}{user1},
  \resourceoptval{sort}{double}\comment{double-precision numeric sort}
}
\end{codebox*}%
The last resource set assumes that the entries defined in the
file \filefmt{constants.bib} have a number stored in the \gloskey{user1}
field. For example:
\begin{compactcodebox}
\atentry{number}\marg{pi,
  \gloskeyval{name}{\cmd{ensuremath}\marg{\cmd{pi}}},
  \gloskeyval{description}{pi},
  \gloskeyval{user1}{3.141592654}
}
\end{compactcodebox}

   \item Put 
\begin{compactcodebox*}
\gls{printunsrtglossary}\oarg{\printglossoptval{type}{\meta{type}},\meta{options}}
\end{compactcodebox*}%
    where you want your list of entries (\idx{glossary}) to appear. (The
    \printglossopt{sort} key isn't available in \meta{options}. The
    \resourceopt{sort} setting needs to be
    used in \gls{GlsXtrLoadResources} instead.)
    Alternatively, use
\begin{compactcodebox*}
\gls{printunsrtglossaries}
\end{compactcodebox*}
    which will display all glossaries (if you have more than one).

    \item Run \LaTeX\ on your document. The \opt{record} option 
    adds information to the \ext+{aux} file that provides
    \app{bib2gls} with all required details for each resource set.
    For example, if the file is called \filefmt{myDoc.tex}:
\begin{terminal}
pdflatex myDoc
\end{terminal}

    \item Run \app{bib2gls}
\begin{terminal}
bib2gls myDoc
\end{terminal}%
or (if you need letter groups)
\begin{terminal}
bib2gls \switch{group} myDoc
\end{terminal}

    \item Run \LaTeX\ again.
   \end{enumerate}
Here's a complete document (\filefmt{myDoc.tex}) with only one
resource set for brevity:
\begin{codebox}
\cmd{documentclass}\marg{article}
\codepar
\cmd{usepackage}[\opt{record}]\marg{glossaries-extra}
\comment{input the \ext{glstex} file created by \app{bib2gls}:}
\gls{GlsXtrLoadResources}
 \oarg{\comment{instructions to \app{bib2gls}:}
   \resourceoptvalm{src}{entries}, \comment{data in entries.bib}
   \resourceoptval{sort}{en-GB}\comment{sort according to this locale}
 }
\codepar
\cbeg{document}
A \gls{gls}\marg{sample}.
\codepar
\comment{iterate over all defined entries:}
\gls{printunsrtglossaries}
\cend{document}
\end{codebox}
The accompanying \filefmt{entries.bib} file:
\begin{compactcodebox}
\atentry{entry}\marg{sample,
  \gloskeyval{name}{sample},
  \gloskeyval{description}{an example}
}
\end{compactcodebox}
Document build:
\begin{terminal}
pdflatex myDoc
bib2gls myDoc
pdflatex myDoc
\end{terminal}
\end{itemdesc}

\end{deflist}

If you are having difficulty integrating \app{makeglossaries} (or
\app{makeglossaries-lite}) or \app{bib2gls}
into your document build process, you may want to consider using
\app{arara}, which is a Java application that searches the
document for special comment lines that tell \app{arara} which
applications to run. For example, the file \filefmt{myDoc.tex} might
start with:
\begin{codebox}
\araraline{pdflatex}
\araraline{makeglossaries}
\araraline{pdflatex}
\cmd{documentclass}\marg{article}
\cmd{usepackage}\marg{glossaries}
\gls{makeglossaries}
\end{codebox}
then to build the document you just need the single system call:
\begin{terminal}
arara myDoc
\end{terminal}
Alternatively, if you want to use \app{makeglossaries-lite}, change
the second line to:
\begin{compactcodebox}
\araraline{makeglossarieslite}
\end{compactcodebox}
There's also a rule for \app{bib2gls}. For example, the following
indicates that letter groups are required:
\begin{codebox}
\araraline{pdflatex}
\araraline{bib2gls: \marg{ group: on }}
\araraline{pdflatex}
\cmd{documentclass}\marg{article}
\cmd{usepackage}\oarg{\opt{record}}\marg{glossaries-extra}
\gls{GlsXtrLoadResources}
\end{codebox}

\section{Customising the Glossary}
\label{sec:glosstyle}

The default glossary style uses the \env{description} environment
to display the entry list. Each entry name is set in the optional
argument of \gls{item} which means that it will typically be
displayed in bold. You can switch to medium weight by redefining
\gls{glsnamefont}:
\begin{codebox}
\cmd{renewcommand}*\marg{\gls{glsnamefont}}[1]\marg{\cmd{textmd}\marg{\#1}}
\end{codebox}
Some classes and packages redefine the \env{description} environment
in such as way that's incompatible with the \sty{glossaries}
package. In which case you'll need to select a different glossary
style (see below).

By default, a~full stop is appended to the description (unless you
use \sty{glossaries-extra}). To prevent
this from happening use the \opt{nopostdot} package option:
\begin{codebox}
\cmd{usepackage}[\opt{nopostdot}]\marg{glossaries}
\end{codebox}
or to add it with \sty{glossaries-extra}:
\begin{codebox}
\cmd{usepackage}[\opt{postdot}]\marg{glossaries-extra}
\end{codebox}

By default, a~location list is displayed for each entry (except
when \gls{printunsrtglossary} is used without \app{bib2gls}). This refers
to the document locations (for example, the page number) where the
entry has been referenced. If you use \optionsor{mkidx,xdy} described in
\sectionref{sec:printglossaries} or \option{b2g} (with \app{bib2gls} and
\sty{glossaries-extra}) then location ranges will be compressed.
For example, if an entry was used on pages~1, 2 and~3, then with
\optionsor{mkidx,xdy,b2g} the location list will appear as 1--3, but
with \option{noidx} it
will appear as 1, 2, 3. If you don't want the locations displayed
you can hide them using the \opt{nonumberlist} package option:
\begin{codebox}
\cmd{usepackage}[\opt{nonumberlist}]\marg{glossaries}
\end{codebox}
or with \app{bib2gls} use \resourceoptval{save-locations}{false} in the
optional argument of the appropriate \gls{GlsXtrLoadResources}
(it's possible to have some resource sets with locations and some
without).

Entries are grouped according to the first letter of
each entry's \gloskey{sort} value. By default a~vertical gap is placed
between letter groups for most of the predefined styles. You can suppress this with the
\opt{nogroupskip} package option:
\begin{codebox}
\cmd{usepackage}[\opt{nogroupskip}]\marg{glossaries}
\end{codebox}

If the default style doesn't suit your document, you can change the
style using:
\cmddef{setglossarystyle}
\dickimawhref{gallery/glossaries-styles/}{There are a~number of predefined styles.} Glossaries can vary from
a~list of symbols with a~terse description to a~list of words or
phrases with descriptions that span multiple paragraphs, so there's
no \qt{one style fits all} solution. You need to choose a~style that
suits your document. For example:
\begin{codebox*}
\gls{setglossarystyle}\marg{\glostyle{index}}
\end{codebox*}
You can also use the \opt{style} package option for the preloaded
styles. For example:
\begin{codebox}
\cmd{usepackage}[\optval{style}{index}]\marg{glossaries}
\end{codebox}

Examples:
\begin{enumerate}
 \item You have entries where the name is a~symbol and the
 description is a~brief phrase or short sentence. Try one of the 
 \qt{mcol} styles defined in the \sty{glossary-mcols} package. For example:
\begin{codebox}
\cmd{usepackage}[\opt{nopostdot}]\marg{glossaries}
\cmd{usepackage}\marg{glossary-mcols}
\gls{setglossarystyle}\marg{\glostyle+{mcolindex}}
\end{codebox}
or
\begin{codebox*}
\cmd{usepackage}[\optval{stylemods}{mcols},\optval{style}{mcolindex}]\marg{glossaries-extra}
\end{codebox*}

 \item You have entries where the name is a~word or phrase and the
description spans multiple paragraphs. Try one of the \qt{altlist} styles. For
example:
\begin{codebox}
\cmd{usepackage}[\opt{nopostdot}]\marg{glossaries}
\gls{setglossarystyle}\marg{\glostyle+{altlist}}
\end{codebox}
or
\begin{codebox}
\cmd{usepackage}[\opt{stylemods},\optval{style}{altlist}]\marg{glossaries-extra}
\end{codebox}

 \item You have entries where the name is a~single word, the
  description is brief, and an associated symbol has been set.
  Use one of the styles that display the symbol (not all of them do). 
  For example, one of the tabular styles:
\begin{codebox}
\cmd{usepackage}[\opt{nopostdot},\opt{nonumberlist},
 \optval{style}{\glostyle+{long4col}}]\marg{glossaries}
\end{codebox}
or one of the \qt{tree} styles:
\begin{codebox}
\cmd{usepackage}[\opt{nopostdot},\opt{nonumberlist},\optval{style}{tree}]\marg{glossaries}
\end{codebox}
\end{enumerate}

If your glossary consists of a~list of abbreviations and you also want to
specify a~description as well as the long form, then you need to use
an abbreviation style that will suit the glossary style. For example,
use the \acrstyle{long-short-desc} acronym style:
\begin{codebox*}
\gls{setacronymstyle}\marg{\acrstyle{long-short-desc}}
\end{codebox*}
Define the acronyms with a~description:
\begin{codebox*}
\gls{newacronym}
 \oarg{\gloskeyval{description}{statistical pattern
   recognition technique}}
 \marg{svm}\marg{SVM}\marg{support vector machine}
\end{codebox*}
Alternatively with \sty{glossaries-extra}:
\begin{codebox*}
\gls{setabbreviationstyle}\marg{\acrstyle{long-short-desc}}
\codepar
\gls{newabbreviation}
 \oarg{\gloskeyval{description}{statistical pattern
    recognition technique}}
 \marg{svm}\marg{SVM}\marg{support vector machine}
\end{codebox*}

Choose a~glossary style that suits wide entry names. For example:
\begin{codebox*}
\gls{setglossarystyle}\marg{\glostyle{altlist}}
\end{codebox*}

\section{Multiple Glossaries}
\label{sec:multigloss}

The \sty{glossaries} package predefines a~default \code{main}
\idx{glossary}. When you define an entry (using one of the commands
described in \sectionref{sec:defterm}), that entry is automatically
assigned to the default \idx{glossary}, unless you indicate otherwise 
using the \gloskey{type} key. However you first need to
make sure the desired \idx{glossary} has been defined. This is done using:
\cmddef{newglossary}
The \meta{glossary-label} is a~label that uniquely identifies this new
\idx{glossary}. As with other types of identifying labels, be careful not
to use active characters in \meta{label}. The final argument
\meta{title} is the section or chapter title used by
\gls{printglossary}, \gls{printnoidxglossary} or
\gls{printunsrtglossary}. The other arguments
indicate the file extensions used by
\app{makeindex} or \app{xindy} (described in
\sectionref{sec:printglossaries}). If you're not using either
\app{makeindex} or \app{xindy}, 
then the \meta{log-ext}, \meta{in-ext} and \meta{out-ext} arguments aren't
relevant, in which case you may prefer to use the starred version
where those arguments are omitted:
\cmddef{newglossary*}

In the case of \optionsor{mkidx,xdy}, all new \idxpl{glossary} must
be defined before \gls{makeglossaries}. (\emph{Entry} definitions
should come after \gls{makeglossaries}.) In the case of
\option{b2g}, all new \idxpl{glossary} must be defined before any
\gls{GlsXtrLoadResources} that requires them.

Since it's quite common for documents to have both a~list of terms
and a~list of abbreviations, the \sty{glossaries} package provides the
package option \opt{acronyms} (or \opt{acronym}), which 
is a~convenient shortcut for
\begin{compactcodebox}
\gls{newglossary}\oarg{\ext+{alg}}\marg{acronym}\marg{\ext+{acr}}\marg{\ext+{acn}}\marg{\gls+{acronymname}}
\end{compactcodebox}
The option also changes the behaviour of \gls{newacronym} so that acronyms
are automatically put in the list of acronyms instead of the main
glossary. The \sty{glossaries-extra} package also supports the
\opt{acronyms} option for abbreviations defined using \gls{newacronym} but
additionally has the package option \opt{abbreviations} to create
a list of abbreviations for \gls{newabbreviation}.

There are some other package options for creating commonly used
lists: \opt{symbols} (lists of symbols), \opt{numbers} (lists
of numbers), \opt{index} (an index). Since indexes don't
typically have descriptions, the \opt{index} option also defines:
\cmddef{newterm}
This is just a shortcut that uses \gls{newglossaryentry} with the
\gloskey{name} set to \meta{entry-label} and the \gloskey{description} is
suppressed.

For example, suppose you want a~main glossary for terms, a~list of
acronyms and a~list of notation:
\begin{codebox}
\cmd{usepackage}[\opt{acronyms}]\marg{glossaries}
\gls{newglossary}\oarg{nlg}\marg{notation}\marg{not}\marg{ntn}\marg{Notation}
\end{codebox}
After \gls{makeglossaries} (or \gls{makenoidxglossaries}) you can
define the entries in the preamble. For example:
\begin{codebox}
\gls{newglossaryentry}\marg{gls:set}
\marg{\comment{This entry goes in the `main' glossary}
  \gloskeyval{name}{set},
  \gloskeyval{description}{A collection of distinct objects}
}
\codepar
\comment{This entry goes in the `acronym' glossary:}
\gls{newacronym}\marg{svm}\marg{svm}\marg{support vector machine}
\codepar
\gls{newglossaryentry}\marg{not:set}
\marg{\comment{This entry goes in the `notation' glossary:}
  \gloskeyval{type}{notation},
  \gloskeyval{name}{\cmd{ensuremath}\marg{\cmd{mathcal}\marg{S}}},
  \gloskeyval{description}{A set},
  \gloskeyval{sort}{S}}
\end{codebox}
or if you don't like using \csfmt{ensuremath}:
\begin{codebox}
\gls{newglossaryentry}\marg{not:set}
\marg{\comment{This entry goes in the `notation' glossary:}
  \gloskeyval{type}{notation},
  \gloskeyval{name}{\$\cmd{mathcal}\marg{S}\$},
  \gloskeyval{text}{\cmd{mathcal}\marg{S}},
  \gloskeyval{description}{A set},
  \gloskeyval{sort}{S}
}
\end{codebox}

Each glossary is displayed using:
\begin{compactcodebox*}
\gls{printnoidxglossary}\oarg{\printglossoptval{type}{\meta{type}}}
\end{compactcodebox*}%
(\option{noidx}) or
\begin{compactcodebox*}
\gls{printglossary}\oarg{\printglossoptval{type}{\meta{type}}}
\end{compactcodebox*}%
(\options{mkidx,xdy}). Where \meta{type} is the glossary label. If the
type is omitted the default \code{main} glossary is assumed.

If you're using \app{bib2gls} then each glossary is displayed
using:
\begin{compactcodebox*}
\gls{printunsrtglossary}\oarg{\printglossoptval{type}{\meta{type}}}
\end{compactcodebox*}%
With this method you don't use \gls{makeglossaries} or
\gls{makenoidxglossaries}. Instead you can assign the entry type
with the resource command. For example:
\begin{codebox*}
\cmd{usepackage}[\opt{record},\opt{abbreviations},\opt{symbols}]\marg{glossaries-extra}
\codepar
\gls{GlsXtrLoadResources}\oarg{
 \resourceoptvalm{src}{terms}, \comment{entries defined in terms.bib}
 \resourceoptval{type}{main}\comment{put in main glossary}
}
\gls{GlsXtrLoadResources}\oarg{
 \resourceoptvalm{src}{abbrvs}, \comment{entries defined in abbrvs.bib}
 \resourceoptval{type}{abbreviations}\comment{put in abbreviations glossary}
}
\gls{GlsXtrLoadResources}\oarg{
 \resourceoptvalm{src}{syms}, \comment{entries defined in syms.bib}
 \resourceoptval{type}{symbols}\comment{put in symbols glossary}
}
\end{codebox*}
Later in the document:
\begin{codebox*}
\gls{printunsrtglossary} \comment{main}
\gls{printunsrtglossary}\oarg{\printglossoptval{type}{abbreviations}}
\gls{printunsrtglossary}\oarg{\printglossoptval{type}{symbols}}
\end{codebox*}

There's a~convenient shortcut that will display all the defined
glossaries depending on the indexing method:
\cmddef{printnoidxglossaries}
(\option{noidx}) or
\cmddef{printglossaries}
(\options{mkidx,xdy}) or (\sty{glossaries-extra} only)
\cmddef{printunsrtglossaries}

If you use \optionsor{noidx,b2g}, you don't need to do anything different
for a document with multiple \idxpl{glossary}. If you use
\optionsor{mkidx,xdy} with the \app{makeglossaries} Perl script or the
\app{makeglossaries-lite} Lua script, you
similarly don't need to do anything different to the document build
(compared to building a document with only one \idx{glossary}).

If you use \optionsor{mkidx,xdy} without the helper \app{makeglossaries}
Perl script or \app{makeglossaries-lite} Lua script then you need to
make sure you run \app{makeindex}\slash\app{xindy} \emph{for each
defined glossary}. The \meta{gls} and \meta{glo} arguments of
\gls{newglossary} specify the file extensions to use instead of
\ext+{gls} and \ext+{glo}. The optional argument \meta{glg} is the
file extension for the transcript file. This should be different for
each glossary in case you need to check for
\app{makeindex}\slash\app{xindy} errors or warnings if things go
wrong.

For example, suppose you have three glossaries in your document
(\code{main}, \code{acronym} and \code{notation}), 
specified using:
\begin{codebox}
\cmd{usepackage}[\opt{acronyms}]\marg{glossaries}
\gls{newglossary}\oarg{nlg}\marg{notation}\marg{not}\marg{ntn}\marg{Notation}
\end{codebox}
Then (assuming your \LaTeX\ document is in a~file called
\filefmt{myDoc.tex}):

\begin{description}
\item\option{mkidx}:
Either use one \app{makeglossaries} or
\app{makeglossaries-lite} call:
\begin{terminal}
makeglossaries myDoc
\end{terminal}
or
\begin{terminal}
makeglossaries-lite myDoc
\end{terminal}
Or you need to run \app{makeindex} three times:
\begin{terminal}
makeindex -t myDoc.glg -s myDoc.ist -o myDoc.gls myDoc.glo
makeindex -t myDoc.alg -s myDoc.ist -o myDoc.acr myDoc.acn
makeindex -t myDoc.nlg -s myDoc.ist -o myDoc.not myDoc.ntn
\end{terminal}

\item\option{xdy}:
Either use one \app{makeglossaries} call:
\begin{terminal}
makeglossaries myDoc
\end{terminal}
Or you need to run \app{xindy} three times (be careful not to insert
line breaks where the line has wrapped.)
\begin{terminal}
xindy -L english -C utf8 -I xindy -M myDoc -t myDoc.glg -o myDoc.gls myDoc.glo
xindy -L english -C utf8 -I xindy -M myDoc -t myDoc.alg -o myDoc.acr myDoc.acn
xindy -L english -C utf8 -I xindy -M myDoc -t myDoc.nlg -o myDoc.not myDoc.ntn
\end{terminal}

\item\option{b2g}:
With \app{bib2gls} only one call is required:
\begin{terminal}
pdflatex myDoc
bib2gls \switch{group} myDoc
pdflatex myDoc
\end{terminal}
(Omit \switch{group} if you don't need letter groups.)
\end{description}

\section{Hyperlinks (\stytext{glossaries} and \stytext{hyperref})}
\label{sec:hyperref}

Take care if you use the \sty{glossaries} package with
\sty{hyperref}. Contrary to the usual advice that \sty{hyperref}
should be loaded last, \sty{glossaries} (and \sty{glossaries-extra}) 
must be loaded \emph{after} \sty{hyperref}:
\begin{codebox}
\cmd{usepackage}[colorlinks]\marg{hyperref}
\cmd{usepackage}\marg{glossaries}
\end{codebox}

By default, if the \sty{hyperref} package has been loaded, commands
like \gls{gls} will form a~hyperlink to the relevant entry in the
\idx{glossary}. If you want to disable this for \emph{all} your
glossaries, then use:
\cmddef{glsdisablehyper}
If you want hyperlinks suppressed for entries in specific
glossaries, then use the \opt{nohypertypes} package option. For
example, if you don't want hyperlinks for entries in the \code{acronym} and
\code{notation} glossaries but you do want them for entries in the
\code{main} glossary, then do:
\begin{codebox}
\cmd{usepackage}[colorlinks]\marg{hyperref}
\cmd{usepackage}[\opt{acronym},\optvalm{nohypertypes}{acronym,notation}]\marg{glossaries}
\gls{newglossary}\oarg{nlg}\marg{notation}\marg{not}\marg{ntn}\marg{Notation}
\end{codebox}

If you want the hyperlinks suppressed the first time an entry is
used, but you want hyperlinks for subsequence references then use
the \optval{hyperfirst}{false} package option:
\begin{codebox}
\cmd{usepackage}[colorlinks]\marg{hyperref}
\cmd{usepackage}[\optval{hyperfirst}{false}]\marg{glossaries}
\end{codebox}
The \sty{glossaries-extra} extension package provides another method
using category attributes. See the \sty{glossaries-extra} user
manual for further details.

Take care not to use non-expandable commands in PDF bookmarks. This
isn't specific to the \sty{glossaries} package but is a~limitation
of PDF bookmarks. Non-expandable commands include commands like 
\gls{gls}, \gls{glspl}, \gls{Gls} and \gls{Glspl}. The \sty{hyperref}
package provides a~way of specifying alternative text for the PDF
bookmarks via \gls{texorpdfstring}. For example:
\begin{codebox}
\cmd{section}\marg{The \gls{texorpdfstring}\marg{\gls{gls}\marg{fishage}}\marg{Fish Age}}
\end{codebox}
However, it's not a~good idea to use commands like \gls{gls} in
a~section heading as you'll end up with the table of contents page in 
your location list and it will unset the \idx{firstuseflag} too soon. 
Instead you can use
\cmddef{glsentrytext}
This is expandable provided that the \gloskey{text} key doesn't
contain non-expandable code. For example, the following works:
\begin{codebox}
\cmd{section}\marg{The \gls{glsentrytext}\marg{fishage}}
\end{codebox}
and it doesn't put the table of contents in the location list.

If you use \sty{glossaries-extra} then use the commands that are
provided specifically for use in section headers. For example:
\begin{codebox*}
\cmd{section}\marg{The \gls{glsfmttext}\marg{fishage}}
\end{codebox*}

\section{Cross-References}
\label{sec:xr}

You can add a~reference to another entry in a~location list using
the \gloskeyval{see}{\meta{label-list}} key when you define an entry. 
The referenced entry (or entries) must also be defined.

For example:
\begin{codebox}
\gls{longnewglossaryentry}\marg{devonian}\marg{\gloskeyval{name}{Devonian}}\comment{}
\marg{\comment{}
  The geologic period spanning from the end of the 
  Silurian Period to the beginning of the
  Carboniferous Period.
\codepar
  This age was known for its remarkable variety of 
  fish species.
}
\codepar
\gls{newglossaryentry}\marg{fishage}
\marg{
  \gloskeyval{name}{Fish Age},
  \gloskeyval{description}{Common name for the Devonian period},
  \gloskeyval{see}{devonian}
}
\end{codebox}
The cross-reference will appear as \qt{\emph{see} Devonian}. You can
change the \qt{see} tag for an individual entry using the format
\gloskeyval{see}{\oargm{tag}\meta{label-list}}. For example:
\begin{codebox}
\gls{newglossaryentry}\marg{latinalph}
\marg{
  \gloskeyval{name}{Latin alphabet},
  \gloskeyval{description}{alphabet consisting of the letters 
  a, \cmd{ldots}, z, A, \cmd{ldots}, Z},
  \gloskey{see}=[see also]\marg{exlatinalph}
}
\gls{newglossaryentry}\marg{exlatinalph}
\marg{
   \gloskeyval{name}{extended Latin alphabet},
   \gloskeyval{description}{The Latin alphabet 
   extended to include other letters such as 
   ligatures or diacritics.}
}
\end{codebox}
In the above, I haven't enclosed the entire value of the
\gloskey{see} key in braces.
If you use the \gloskey{see} key in an optional argument, such as
the optional argument of \gls{newacronym}, make sure you enclose the 
value (including the optional tag) in braces. For example:
\begin{codebox}
\gls{newacronym}\marg{ksvm}\marg{ksvm}\marg{kernel support vector machine}
\gls{newacronym}
 \oarg{\gloskeyval{see}{[see also]\marg{ksvm}}}
 \marg{svm}\marg{svm}\marg{support vector machine}
\end{codebox}
The \sty{glossaries-extra} package provides a \gloskey{seealso}
key. This doesn't allow a tag but behaves much like
\gloskeyval{see}{[\gls{seealsoname}]\margm{label-list}}. For example:
\begin{codebox}
\gls{newabbreviation}\marg{ksvm}\marg{ksvm}
  \marg{kernel support vector machine}
\gls{newabbreviation}
 \oarg{\gloskeyval{seealso}{ksvm}}
 \marg{svm}\marg{svm}\marg{support vector machine}
\end{codebox}
Additionally, the \sty{glossaries-extra} package provides
an \gloskey{alias} key where the value must be just a single label.
See \gallerypage{aliases}{glossaries-extra gallery: aliases} for an example.

Since the cross-reference appears in the location list, if you
suppress the location list using the \opt{nonumberlist} package
option, then the cross-reference will also be suppressed. With
\app{bib2gls}, don't use the \opt{nonumberlist} package
option in this case, but instead use 
\resourceoptval{save-locations}{false} in the resource
options. For example:
\begin{codebox}
\cmd{usepackage}[\opt{record},\opt{abbreviations},\opt{symbols}]\marg{glossaries-extra}
\codepar
\gls{GlsXtrLoadResources}\oarg{
 \resourceoptvalm{src}{terms}, \comment{entries defined in terms.bib}
 \resourceoptval{type}{main}\comment{put in main glossary}
}
\gls{GlsXtrLoadResources}\oarg{
 \resourceoptvalm{src}{abbrvs}, \comment{entries defined in abbrvs.bib}
 \resourceoptval{type}{abbreviations},\comment{put in abbreviations glossary}
\comment{no number list for these entries:}
 \resourceoptval{save-locations}{false}
}
\gls{GlsXtrLoadResources}\oarg{
 \resourceoptvalm{src}{syms}, \comment{entries defined in syms.bib}
 \resourceoptval{type}{symbols}\comment{put in symbols glossary}
}
\end{codebox}

\section{Further Information}
\label{sec:moreinfo}

\begin{itemize}
\item \ctanmirrordocnofn{support/bib2gls/bib2gls-begin}{\styfmt{glossaries-extra}
and \appfmt{bib2gls}: an introductory guide}.
\item The main \sty{glossaries} \docref{user manual}{glossaries-user}.
\item The \dickimawhrefnofn{faqs/glossariesfaq.html}{glossaries FAQ}.
\item \dickimawhrefnofn{latex/buildglossaries/}{Incorporating
makeglossaries or makeglossaries-lite or bib2gls into the document
build}.
\item The \href{https://ctan.org/pkg/glossaries-extra}{\styfmt{glossaries-extra}} package.
\item The \href{https://ctan.org/pkg/bib2gls}{\appfmt{bib2gls}}
application.
\end{itemize}
The \gallery{Dickimaw Books Gallery} provides additional example documents.

\printterms
\printsummary
\clearpage
\printuserguideindex

\end{document}
